% !TEX program = xelatex
\documentclass[12pt, a4paper]{article}

\usepackage{fontspec}
\setmainfont{TeX Gyre Pagella}
\usepackage{geometry}
\geometry{margin=1.2in}
\usepackage{setspace}
\onehalfspacing
\usepackage{amsmath, amssymb, amsthm}
\usepackage{mathtools}
\usepackage{hyperref}
\hypersetup{colorlinks=true, linkcolor=black, citecolor=black, urlcolor=black}
\usepackage{titlesec}
\usepackage[round]{natbib}

\titleformat{\section}{\large\bfseries}{}{0em}{}[\titlerule]
\titleformat{\subsection}{\normalsize\bfseries}{}{0em}{}

\newtheoremstyle{essaystyle}
  {1em}{1em}{\itshape}{}{\bfseries}{.}{0.5em}{}
\theoremstyle{essaystyle}
\newtheorem{definition}{Definition}
\newtheorem{proposition}{Proposition}
\newtheorem*{remark}{Remark}

\title{\textbf{Function Survives Collapse}\\[0.4em]
\large Admissibility, Entropy Flow, and the Geometry\\of Persistent Computation}
\author{Flyxion\\[0.3em]
\small Independent Researcher}
\date{May 2026}

\begin{document}

\maketitle
\thispagestyle{empty}

\begin{abstract}
Three results from distinct scientific domains converge on a single structural
principle. In cryptography, the theory of additive randomized encodings demonstrates
that arbitrary secure computation is achievable through anonymous additive aggregation:
local provenance is irreversibly destroyed while computable functional invariants persist.
In molecular biology, a study of transcription--translation coupling in
\textit{Escherichia coli} and related organisms demonstrates that gene expression
is not a linear pipeline but a geometry-constrained coordination of partially decoupled
flows, where mRNA fate depends critically on spatial accessibility rather than
informational sequence alone. In plant developmental biology, a single-nucleus
transcriptional atlas of early \textit{Arabidopsis} seed development reveals that
the embryo, endosperm, and seed coat coordinate their developmental programs not
through full informational transparency but through low-dimensional signalling
projections across symplastically isolated tissue interfaces. This essay argues
that all three results are instances of a single deeper principle, formalized through
the notion of admissible projection: a map $\pi : X \to M$ from a microscopic
trajectory space $X$ onto an admissible invariant manifold $M$, which preserves
computable functionals while sacrificing microscopic recoverability under constrained
entropy flow. The Relativistic Scalar-Vector Plenum framework supplies the overarching
field-theoretic interpretation, with a variational Lagrangian formulation deriving the
admissible projection as the attractor of constrained Euler--Lagrange flow. The
framework is further extended through a categorical treatment of morphisms between
admissible systems, a renormalization-theoretic interpretation of scale-dependent
projection, and a formal analysis of catastrophic projection collapse as the failure
mode of robust systems. The central claim: successful computational, biological, and
developmental systems do not preserve microscopic structure. They preserve
low-dimensional functional invariants while continuously destroying the recoverability
of the paths by which those invariants were reached.
\end{abstract}

\tableofcontents
\newpage

%------------------------------------------------------
\section{Introduction: Two Systems, One Collapse}
%------------------------------------------------------

Consider two facts from different sciences, presented without interpretation.

First: if $n$ parties wish to compute a joint function of their private inputs without
revealing anything beyond the output, they can each independently encode their input
as a randomized group element and submit all encodings to an anonymous shuffle.
The shuffler sums the encodings. The evaluator applies a decoder. The function output
is recovered. No individual input is recoverable. Not approximately concealed ---
statistically indistinguishable from a transcript generated knowing only the output.
This is the additive randomized encoding (ARE) framework. The foundational framework
and its initial feasibility results appear in \citet{halevi2023additive}; the full
universality result, establishing that statistical ARE exist for all functions, is due
to \citet{shufflinguniversal}, who refuted the earlier conjecture that statistical
security fails for simple functions. Efficient constructions remain available only for
restricted circuit classes.

Second: in \textit{Escherichia coli}, the enzyme RNase~E is the primary driver of
messenger RNA degradation. RNase~E is not distributed uniformly through the cell.
It is anchored to the inner membrane. Transcription occurs deep in the nucleoid.
The geometric separation between these two regions means that nascent transcripts,
still tethered to the RNA polymerase and threaded through the translation apparatus,
are spatially inaccessible to the degradation machinery. Co-transcriptional degradation
is therefore negligible --- not because of an active protective mechanism, but because
the geometry forbids the encounter. When RNase~E is experimentally relocated to the
cytoplasm, co-transcriptional degradation rises approximately sevenfold
\citep{kim2026spatial}.

These phenomena appear to belong to incommensurable domains. One concerns
privacy-preserving distributed computation. The other concerns the subcellular
organization of bacterial gene regulation. Their vocabularies, methods, and epistemic
standards have essentially no overlap.

Yet both are instances of the same structural situation: a system in which computable
functional output persists under conditions that make microscopic reconstruction
impossible. In the ARE case, impossibility is the design criterion. In the biological
case, impossibility is a consequence of geometry. But in both cases the invariant that
survives is low-dimensional relative to the space of microscopic configurations that
could in principle produce it, and the mechanism of persistence is not protection of the
path but indifference to it.

This convergence is not coincidental, and it is not merely analogical. Both systems
instantiate the same abstract structure: an irreversible projection from a
high-dimensional microscopic trajectory space onto a low-dimensional admissible
invariant manifold. The remainder of this essay develops that structure formally,
applies it to each domain in detail, and argues that the Relativistic Scalar-Vector
Plenum (RSVP) framework --- a field-theoretic account of constrained dynamical systems
developed for cosmological, cognitive, and computational applications --- provides the
natural overarching geometry for the unification.

The central thesis, stated once and carried through what follows: both systems preserve
low-dimensional functional invariants by sacrificing microscopic recoverability under
constrained entropy flow. The functional output is more stable than the path to it.

\subsection*{Structure of the Essay}

Part~I develops the core formalism: the admissible projection map, the notion of
invariant persistence, and the role of entropy flow in destroying provenance while
preserving output. Part~II applies this to the ARE and shuffle-model MPC framework,
showing how universal computation emerges from entropy mixing over locally randomized
fragments. Part~III applies it to transcription--translation coupling and mRNA
degradation in bacteria, showing how spatial admissibility constraints govern which
molecular trajectories can interact. Part~IV develops the unified RSVP interpretation
and introduces the field triple $(\Phi, \mathbf{v}, S)$ as its formal completion.
Part~V develops the generalization of privacy as accessibility geometry. Part~VI
examines constraint geometry, entropy as functional selector, and persistent computation
without persistent representation. Part~VII extends the framework toward a general
theory of admissible systems and addresses the philosophical implications for microscopic
ontology. The conclusion reflects on what it means for function to survive collapse.

%------------------------------------------------------
\section{Part I: Admissible Projection and Invariant Persistence}
%------------------------------------------------------

\subsection{The Core Formalism}

Let $X$ denote a space of microscopic configurations or trajectories. In the
cryptographic setting, $X$ is the space of all possible joint inputs and randomness
choices across $n$ parties. In the biological setting, $X$ is the space of all possible
molecular trajectory ensembles: positions, velocities, binding states, and temporal
sequences of every relevant molecular species. In both cases, $X$ is vast,
high-dimensional, and inaccessible to direct observation.

\begin{definition}[Admissible Invariant Manifold]
Let $\mathcal{F} = \{f_\alpha\}$ be a collection of computable functionals
$f_\alpha : X \to Y_\alpha$ representing the observables of interest. The
\emph{admissible invariant manifold} $M$ is the subspace of $X$, or of a suitable
quotient thereof, on which the values of all $f_\alpha$ are well-defined and stable
under small perturbations consistent with the constraints of the system. Formally,
\[
M = \bigl\{ [x] \in X/{\sim} \;\big|\; f_\alpha(x) = f_\alpha(x') \text{ for all }
x \sim x',\; \forall \alpha \bigr\},
\]
where $x \sim x'$ if $x$ and $x'$ are indistinguishable under the constraint
structure of the system.
\end{definition}

\begin{definition}[Admissible Projection]
An \emph{admissible projection} is a map
\[
\pi : X \longrightarrow M
\]
satisfying four conditions. First, functional preservation: $f_\alpha(\pi(x)) =
f_\alpha(x)$ for all $x \in X$ and all $\alpha$. Second, irreversibility: $\pi$ is not
injective; multiple distinct elements of $X$ map to the same element of $M$. Third,
entropy non-decrease: $H(\pi(X)) \leq H(X)$ in the sense that the projection increases
or preserves the entropy of the induced distribution. Fourth, provenance destruction:
there exists no efficient procedure for recovering $x$ from $\pi(x)$ and the constraint
specification of the system.
\end{definition}

\begin{remark}
The first condition ensures that the projection is not merely information loss: it is
structure-preserving loss. The system discards everything about $x$ except what is
needed to compute the designated outputs. The second and fourth conditions together
formalize what it means for microscopic recoverability to fail. The third condition
locates this failure within a thermodynamic frame. The projection is not a filter that
removes noise while preserving signal. It is a collapse that treats the microscopic
trajectory as signal and the functional output as the invariant to be preserved.
\end{remark}

\subsection{Entropy Flow and the Collapse of Provenance}

The language of entropy flow is necessary here because neither paper under examination
treats its subject as static. Both describe dynamic processes in which configurations
evolve, interact, and are partially destroyed before a stable output is reached.

In the ARE setting, the relevant entropy flow is the injection of randomness by local
encoders followed by aggregation. Each party $P_i$ with input $x_i$ samples randomness
$r_i$ and produces an encoding $\hat{x}_i = \mathrm{Enc}(x_i; r_i)$. The evaluator
receives $\hat{y} = \sum_i \hat{x}_i$. The randomness injected by each encoder
thermalizes the local representation: $\hat{x}_i$ carries no more information about
$x_i$ than is strictly necessary to reconstruct $f(x_1, \ldots, x_n)$ from the
aggregate.

In the biological setting, the relevant entropy flow is the diffusion, translation, and
partial degradation of RNA transcripts across an inhomogeneous spatial field. The
transcript does not persist as a stable object. It is a dynamical entity whose fate
depends on which molecular encounters it undergoes, in which order, at which locations.
The aggregate observable --- protein production rate, mRNA lifetime, coupling ratio ---
is not a property of the transcript in isolation. It is a functional of the trajectory
ensemble under the constraints imposed by the geometry of the cell.

In both cases, the functional output is more stable than the path to it. This is the
operational content of admissible projection: $M$ is an attractor, not an artifact.
The trajectory is consumed. The invariant persists.

\subsection{Pipeline Ontologies and Their Failure}

Both domains under examination were historically understood through pipeline models.
In theoretical computer science and distributed systems, the classical model of secure
computation assumes that parties exchange messages through authenticated channels,
that computations proceed in rounds, and that security is a property of the protocol
architecture. The ARE framework challenges this picture fundamentally: it asks what can
be computed when the only permitted global operation is addition --- no rounds, no
authenticated channels, no correlated setup. The answer, recently established, is:
everything, in the statistical setting. The pipeline from sender to receiver is
unnecessary. Coordinated entropy injection suffices.

In molecular biology, the classical Central Dogma --- DNA $\to$ RNA $\to$ protein ---
carries implicit pipeline semantics: transcription produces a transcript, translation
consumes it, the two processes are temporally and functionally sequential. The paper by
Kim et al.\ \citep{kim2026spatial} systematically dismantles this picture.
Transcription and translation are not obligatorily coupled. Their degree of coupling
depends on translation initiation rate. mRNA degradation does not simply follow
transcript release; it depends on spatial accessibility to the degradation apparatus,
which is itself a function of subcellular geometry and molecular identity. The pipeline
from DNA to protein is not a pipeline. It is a partially coupled dynamical system whose
observable outputs depend on constraint geometry rather than sequential information flow.

The admissible projection framework offers a replacement for the pipeline ontology.
Instead of asking how information travels from source to output, it asks which
functionals of the source are preserved under the constraint-governed collapse of
microscopic recoverability. That reorientation is the conceptual core of both results
under examination.

%------------------------------------------------------
\section{Part II: Additive Randomized Encodings as Entropy Quotient Computation}
%------------------------------------------------------

\subsection{The ARE Framework}

The ARE model \citep{halevi2023additive} formalizes the following question: given an
$n$-party function $f(x_1, \ldots, x_n)$ and an Abelian group $G$, can each party
$P_i$ independently encode its input $x_i$ into a group element $\hat{x}_i \in G$,
such that the sum $\hat{y} = \sum_i \hat{x}_i$ allows recovery of $f(x_1, \ldots, x_n)$
while revealing nothing else? A triple $\Pi = (\mathrm{Setup}, \mathrm{Enc},
\mathrm{Dec})$ constitutes an ARE scheme if there exists a simulator $\mathrm{Sim}$
such that the distribution of $(pp, \hat{y})$ produced by the encoding procedure is
indistinguishable from the distribution produced by $\mathrm{Sim}(f(x_1, \ldots, x_n))$
alone. The evaluator's view is simulatable from the output. No additional information
is extractable.

\begin{definition}[ARE Security]
An ARE scheme $\Pi$ for $f$ is \emph{statistically secure} if there exists a simulator
$\mathrm{Sim}$ such that for all $x_1, \ldots, x_n$,
\[
\mathrm{SD}\bigl(\mathrm{Sim}(f(x_1,\ldots,x_n)),\; \Pi(x_1,\ldots,x_n)\bigr) \leq \delta(\lambda),
\]
where $\delta$ is negligible and $\mathrm{SD}(\cdot,\cdot)$ denotes statistical
distance. The scheme is \emph{perfectly secure} if $\delta = 0$.
\end{definition}

In the language of Part~I, the encoding map $x_i \mapsto \hat{x}_i$ is a local section
of the admissible projection. The summation operator is the aggregation mechanism. The
resulting $\hat{y}$ lives in $M$: it is the image of $\pi(x_1, \ldots, x_n)$,
containing exactly the information needed to decode $f$ and nothing more.

\subsection{The OR Construction as Phase Transition}

The simplest nontrivial ARE is the scheme for the $\mathrm{OR}$ function over $n$
binary inputs \citep{halevi2023additive}. Each party encodes $x_i = 0$ as the identity
element $0 \in G$ and $x_i = 1$ as a uniformly random element $g_i \in_R G$. The sum
$\hat{y} = \sum_{i=1}^n \hat{x}_i$ is zero when $\mathrm{OR}(x) = 0$ and uniformly
random in $G$ (with probability $1 - |G|^{-1}$) when $\mathrm{OR}(x) = 1$.

The evaluator decodes: output $0$ if $\hat{y} = 0$, output $1$ otherwise. Security is
perfect when $\mathrm{OR}(x) = 1$: the transcript $\hat{y}$ is indistinguishable from
random and carries no information about which party contributed the active input, how
many parties were active, or what the individual inputs were.

This construction instantiates a phase transition in the aggregate state. The presence
of a single active participant thermalizes the sum. One unit of entropy injection is
sufficient to mask the entire microscopic configuration. The aggregate observable shifts
from a stable fixed point --- the identity element --- to a uniformly distributed random
variable. The trajectory space $X$ for OR contains $2^n$ configurations. The admissible
manifold $M$ has two elements: the identity and the uniform distribution over $G$.
The projection $\pi : X \to M$ maps all configurations with any active party to the
same distributional output. Provenance is destroyed by the projection. The functional
output is preserved exactly.

\subsection{Universality Through Entropy Amplification}

The deeper result, established by \citet{shufflinguniversal} and refuting the earlier
conjecture that statistical ARE might fail even for simple functions like equality over
small domains, is that statistically secure AREs exist for all functions, without
efficiency guarantees for general functions.

The construction proceeds in two conceptually distinct stages. First, a \emph{leaky}
ARE (LARE) is constructed for any finite function. A LARE satisfies a relaxed security
condition governed by a specific leakage function:
\[
\ell(x, y) =
\begin{cases}
x & \text{if } f(x,y) = 1, \\
\perp & \text{if } f(x,y) = 0.
\end{cases}
\]
That is, when the output is 1, the encoding may reveal the first party's input $x$
while concealing the second party's input $y$ entirely; when the output is 0, nothing
leaks beyond the output. The LARE construction encodes party~1's input $x$ as a
randomized vector with a zero in coordinate $x$ and random group elements elsewhere,
and party~2's input $y$ as a randomized vector with zeros precisely at coordinates
where $f(\cdot, y) = 1$. The aggregate sum therefore has a zero in position $x$ if
and only if $f(x,y) = 1$, yielding correct decoding and the specified leakage profile.
This is perfectly leaky-secure and correct up to a $O(|G|^{-1})$ error.

Second, the LARE is lifted to a full ARE through secret sharing and amplification.
For a security parameter $\tau$, party~1 samples $\tau$ additive secret shares
$\mathbf{x}_1, \ldots, \mathbf{x}_\tau$ of the unit vector $e_x$, and party~2
samples $\tau$ shares $r_1, \ldots, r_\tau$ of $0$. The parties run $\tau$ parallel
LARE instances for the inner-product-with-mask function
$g(\mathbf{a}, (\mathbf{b}, r)) = \langle \mathbf{a}, \mathbf{b}\rangle + r$,
obtaining outputs $\langle \mathbf{x}_i, f_y\rangle + r_i$ for each $i$. Their sum
recovers $\langle e_x, f_y\rangle + 0 = f(x,y)$. Security follows because with
probability $1 - 2^{-\tau+1}$ at least one share $r_i$ ensures the corresponding
output is 0, at which coordinate nothing leaks; whenever the output is 1, the leaked
share $\mathbf{x}_i$ is uniformly random and independent of the actual input $x$.

For any two-party function $f : D_1 \times D_2 \to \{0,1\}$, the reduction to inner
product is achieved by representing party~1's input $x$ as the unit vector $e_x$ and
party~2's input $y$ as the truth-table vector $f_y$ encoding all values $f(\cdot, y)$:
\[
f(x, y) = \langle e_x,\; f_y \rangle.
\]
Function evaluation becomes geometric projection. Computation reduces to inner product,
and inner product reduces to the LARE gadget. Extension to multi-party functions
proceeds via decomposable randomized encodings and oblivious transfer, yielding
efficient constructions for $NL/\mathrm{poly}$ and, assuming one-way functions,
for all of $P/\mathrm{poly}$.

\begin{proposition}[\citealt{shufflinguniversal}, Theorem~1.1]
For any $k$-party function $f : D^k \to D$ over a finite domain and any error bound
$\varepsilon > 0$, there exists an ARE for $f$ with statistical correctness and
security errors at most $\varepsilon$. Moreover, any function $f \in NL/\mathrm{poly}$
admits a polynomial-time efficient ARE with statistical correctness and security errors
$2^{-\tau}$, and any function $f \in P/\mathrm{poly}$ admits a computationally secure
ARE of size $O(\tau \cdot S)$ where $S$ is the circuit size of $f$, using only a
one-way function.
\end{proposition}

The efficiency cost is real and large. But the existence result is philosophically
decisive: the only global operation required is addition. All coordination is
accomplished through local entropy injection and the accumulation of admissible
invariant structure in the aggregate. The pipeline from sender to receiver does not
merely disappear; it is shown to have been unnecessary.

\subsection{The Shuffle as Entropy Quotient Map}

In the shuffle model of distributed computation, the architectural primitive is not a
channel between parties but a mechanism that receives all messages and delivers them in
a uniformly random permutation, erasing sender identity. The ARE framework interfaces
with the shuffle model as follows: the ARE messages $\hat{x}_i$ are submitted to the
shuffle, which sums them while destroying all record of which message came from which
sender \citep{halevi2023additive}.

This shuffler is naturally interpreted as an entropy quotient map. Let $\Sigma$ be the
set of all sender-labeled message multisets. The shuffle defines an equivalence relation
on $\Sigma$: two labeled multisets are equivalent if they produce the same unlabeled
aggregate. The quotient $\Sigma/{\sim}$ is the image of the shuffle. The map
$\sigma : \Sigma \to \Sigma/{\sim}$ destroys provenance --- the sender labels ---
while preserving the additive invariant, the sum.

In the language of Part~I, $\sigma$ is an instance of admissible projection. The fiber
$\pi^{-1}(\hat{y})$ over any aggregate value $\hat{y}$ contains all possible sender
assignments consistent with that aggregate. These fibers are arbitrarily large; the
evaluator cannot determine which fiber element was the actual configuration. The
functional output, recovered by applying $\mathrm{Dec}$ to $\hat{y}$, is
fiber-invariant: it does not depend on which element of $\pi^{-1}(\hat{y})$ was
realized.

\begin{remark}
The cryptographic security condition is equivalent to the geometric statement that the
functional observable is constant on fibers of $\pi$. This is the core content of
simulation-based security, reexpressed in the language of admissible projection. No
simulator is required to reconstruct the fiber element, the actual inputs. It need only
reproduce the fiber-invariant observable, the output distribution.
\end{remark}

\subsection{Robust ARE and the Geometry of Residual Leakage}

The 2023 paper introduces a strengthened notion: \emph{robust ARE}, in which security
holds not only against an external evaluator but against coalitions of parties who may
pool their encodings to learn the partial sum contributed by the remaining honest
parties. In robust ARE, the security guarantee is that the partial sum
$\hat{y}_H = \sum_{i \in H} \hat{x}_i$ reveals no more than the \emph{residual
function} $f_{H,x}$: the function $f$ with the inputs of all honest parties fixed.
Colluding parties can compute all completions of the function, but they cannot learn
the individual inputs of the honest parties beyond what the residual function forces.

This notion has a natural interpretation in obstruction geometry. The information the
adversary cannot extract corresponds to a class recording the failure of local-to-global
extension. The colluding parties have local information --- their own inputs and the
partial sum --- but cannot extend it to a global reconstruction of the honest parties'
inputs. The obstruction to this extension is precisely the entropy injected by the
honest parties' local encoders. The topology of the honest parties' encoding space
prevents the adversary from reconstructing a global section from local data.

%------------------------------------------------------
\section{Part III: Transcription, Coupling, and Spatial Admissibility}
%------------------------------------------------------

\subsection{The Geometry of Bacterial Gene Expression}

The paper by Kim et al.\ \citep{kim2026spatial} investigates the degree to which
transcription, translation, and mRNA degradation are co-temporally coupled in bacteria.
The received view holds that these three processes are tightly integrated: the ribosome
follows closely behind the RNA polymerase, translating the transcript as it emerges,
while the degradation machinery stands ready to process released transcripts promptly.

The paper systematically undermines this picture. Its central finding is that
transcription--translation coupling is not a universal feature of bacterial gene
expression but a condition-dependent outcome governed by translation initiation rate.
Genes with weak ribosome binding sites exhibit significant premature transcription
termination --- the RNA polymerase terminates early because no ribosome is close enough
to prevent it --- while genes with strong ribosome binding sites maintain tight
coupling. The progression factor, defined as the ratio of downstream to upstream
transcript coverage, correlates strongly with translation efficiency genome-wide,
with the correlation coefficient $r = -0.93$ for the engineered variant panel and
a genome-wide separation of high-PF and low-PF gene classes upon Rho inhibition.

This is a quantitative demonstration that the pipeline model of bacterial gene
expression is wrong. There is no universal coupling. There is a coupling field whose
local value depends on the kinetic accessibility of the ribosome to the nascent
transcript.

\subsection{Membrane Localization as Topological Obstruction}

The most striking single result in the paper, from the perspective of admissible
projection, is the role of RNase~E localization. RNase~E, the primary endoribonuclease
responsible for mRNA degradation in \textit{E.~coli}, is anchored to the inner membrane
via its membrane-targeting sequence. Transcription occurs in the nucleoid, which is
spatially separated from the inner membrane. This geometric separation creates a
\emph{spatial admissibility constraint}: the trajectories of nascent transcripts and the
trajectories of membrane-bound RNase~E are largely non-intersecting.

Let $\Gamma_{\mathrm{RNase}}$ denote the spatial trajectory ensemble of RNase~E
molecules and $\Gamma_{\mathrm{nascent}}$ denote the spatial trajectory ensemble of
nascent transcripts still tethered to the RNA polymerase in the nucleoid. The paper's
central finding can be stated geometrically as
\[
\Gamma_{\mathrm{RNase}} \cap \Gamma_{\mathrm{nascent}} \approx \varnothing.
\]
This intersection emptiness is not an active protective mechanism. No molecular
machinery enforces the separation. The geometry of the cell --- the localization of the
membrane, the compaction of the nucleoid, the diffusion rates of the relevant molecules
--- produces the separation as a consequence of spatial organization.

The experimental confirmation is decisive. When RNase~E is relocated to the cytoplasm
by deleting its membrane-targeting sequence, the co-transcriptional degradation rate
$k_{d1}$ rises from $0.042 \pm 0.060 \ \mathrm{min}^{-1}$ to $0.31 \pm 0.08 \
\mathrm{min}^{-1}$, an approximately sevenfold increase \citep{kim2026spatial}. The
functional invariant --- negligible co-transcriptional degradation --- was maintained
not by any biochemical protection mechanism but by the topological separation of
admissible trajectory spaces.

\begin{definition}[Spatial Admissibility Constraint]
A \emph{spatial admissibility constraint} on a molecular interaction is a condition of
the form
\[
\Gamma_A \cap \Gamma_B = \varnothing \pmod{\epsilon},
\]
for molecular species $A$ and $B$, where $\Gamma_A$ and $\Gamma_B$ are their respective
spatial trajectory ensembles and $\epsilon$ is a threshold reflecting diffusion
timescales relative to the dynamical timescale of the interaction being considered. The
constraint renders the interaction kinetically inaccessible regardless of thermodynamic
favorability.
\end{definition}

\subsection{Translation Efficiency as Coupling Field}

The genome-wide analysis in \citet{kim2026spatial} establishes that translation
efficiency --- a proxy for translation initiation rate per transcript --- is the primary
determinant of transcription--translation coupling across the \textit{E.~coli} genome.
Among the 1,336 genes analyzed, high-TE genes show progression factors near 1,
reflecting full coupling, while progressively lower TE values associate with
progressively lower progression factors, corresponding to increased premature
termination.

This finding transforms TE from a purely informational quantity --- a rate of protein
production --- into a field variable governing the dynamical coherence of the
transcription--translation system. High TE maintains the RNAP--ribosome distance below
the threshold of approximately 80 nucleotides at which Rho factor loading and premature
termination become probable. Low TE allows the RNAP to outpace the ribosome, opening
the spacer, enabling Rho loading, and driving the system into a premature termination
regime.

In the admissible projection language, TE parametrizes the projection $\pi$. High-TE
genes are projected onto an admissible manifold $M_{\mathrm{coupled}}$ in which
full-length transcript production and translation are coordinated. Low-TE genes are
projected onto $M_{\mathrm{uncoupled}}$, a different region of the invariant manifold
characterized by partial transcript release, altered degradation kinetics, and two-phase
mRNA decay signatures. The transition between these regimes is continuous rather than
sharp: TE acts as a tunable coupling constant for the field governing spatial and
temporal coordination between the molecular flows.

\subsection{Premature Termination as Trajectory Collapse}

The phenomenon of premature transcription termination is especially illuminating through
the lens of admissible projection. In a strongly coupled gene, the RNAP--ribosome
complex produces a coherent joint trajectory. The transcript is never freely released
into the cytoplasm during transcription; it remains tethered to both the polymerase and
the ribosome. This coupled trajectory is protected from premature termination and, as
argued above, from co-transcriptional degradation by spatial separation from RNase~E.

In a weakly coupled gene, the RNAP outpaces the ribosome. The spacer RNA becomes
accessible to the Rho helicase. Rho loads onto the exposed transcript and drives
termination. The nascent transcript is released prematurely --- incomplete, lacking
the 3' UTR or intrinsic terminator stem-loop structure, and now diffusing freely in
the cytoplasm.

This released transcript occupies a different region of $X$. Its spatial trajectory
ensemble has changed from $\Gamma_{\mathrm{nascent}}$, tethered and nucleoid-proximal,
to $\Gamma_{\mathrm{free}}$, diffusing and cytoplasm-accessible. Critically,
$\Gamma_{\mathrm{free}}$ is no longer disjoint from $\Gamma_{\mathrm{RNase}}$. The
premature transcript becomes accessible to membrane-localized RNase~E through normal
diffusion and is rapidly degraded. The apparent increase in the apparent co-transcriptional
degradation rate $k_{d1}$ observed in weak-RBS strains is therefore not genuinely
co-transcriptional. The paper carefully distinguishes $k_{d1}^*$ --- true
co-transcriptional degradation of RNAP-tethered transcripts --- from $k_{d1}$,
the apparent degradation rate during the transcription window, which includes decay of
prematurely released transcripts. The correlation $r = 0.94$ between premature
termination probability and $k_{d1}$ confirms that the elevated apparent rate is a
consequence of trajectory collapse: the transcript exits the admissible coupled
trajectory and enters a new trajectory ensemble with different accessibility and
degradation kinetics.

\subsection{Species-Specific Admissibility Topologies}

One of the philosophically most significant aspects of \citet{kim2026spatial} is its
comparative analysis across three bacterial species with fundamentally different
subcellular organizations.

In \textit{E.~coli} and \textit{Bacillus subtilis}, the primary ribonuclease ---
RNase~E and its functional homologue RNase~Y, respectively --- is membrane-localized.
In both species, co-transcriptional degradation is negligible for standard cytoplasmic
protein-encoding genes: $k_{d1} \approx 0.042 \ \mathrm{min}^{-1}$ for
\textit{E.~coli} lacZ and $k_{d1} \approx 0.025 \ \mathrm{min}^{-1}$ for
\textit{B.~subtilis} lacZ, both far below the corresponding post-transcriptional rates.
The spatial admissibility constraint $\Gamma_{\mathrm{RNase}} \cap
\Gamma_{\mathrm{nascent}} \approx \varnothing$ holds in both organisms as a consequence
of membrane localization.

In \textit{Caulobacter crescentus}, the situation is fundamentally different. RNase~E
in this organism and its relatives is cytoplasmic. The paper observes rapid
co-transcriptional degradation in \textit{C.~crescentus}: $k_{d1} \approx 1.3 \
\mathrm{min}^{-1}$, approximately an order of magnitude higher than in
\textit{E.~coli}. The spatial admissibility constraint has collapsed.
$\Gamma_{\mathrm{RNase}}$ now overlaps substantially with $\Gamma_{\mathrm{nascent}}$,
making co-transcriptional degradation accessible and prevalent.

\begin{remark}
This comparison establishes that the admissibility topology is not a universal feature
of bacterial gene expression but a species-specific geometric property determined by the
subcellular localization of the primary degradation machinery. The same functional
invariant --- controlled gene expression --- survives in all three organisms. The
microscopic mechanism by which it is achieved differs entirely. Identical informational
goals are realized through distinct geometric accessibility structures. This is better
understood not merely as different evolutionary solutions to the same problem, but as
distinct admissibility geometries realizing the same invariant class.
\end{remark}

The three species thus instantiate three distinct admissible projection geometries for
the same functional task. This is not evolution finding three paths to the same
solution. It is three different implementations of the same abstract admissible
projection $\pi : X \to M$, with different constraint structures on $X$ producing
different intermediate trajectory ensembles while converging on the same class of
invariant observables.

%------------------------------------------------------
\section{Part IV: The RSVP Framework and the Unified Geometry}
%------------------------------------------------------

\subsection{The Field Triple}

The Relativistic Scalar-Vector Plenum (RSVP) framework posits that the relevant
structure of any sufficiently complex dynamical system is captured by a field triple
$(\Phi, \mathbf{v}, S)$, where $\Phi : \mathcal{D} \to \mathbb{R}$ is a scalar field
representing the density or intensity of admissible configurations at each point of the
domain $\mathcal{D}$; $\mathbf{v} : \mathcal{D} \to T\mathcal{D}$ is a vector field
representing the flow of trajectories across the domain; and $S : \mathcal{D} \to
\mathbb{R}$ is an entropy field representing the local degree of microscopic
indistinguishability or randomization.

The dynamics of RSVP systems are governed by the interplay of these three fields.
Regions of high $\Phi$ are regions of high admissibility: trajectories that reach these
regions are likely to produce stable functional outputs. Regions of high $S$ are regions
of high entropy: microscopic recoverability is low, and the projection $\pi$ is far
from injective. The vector field $\mathbf{v}$ governs how trajectories flow from
high-recoverability regions toward high-entropy, high-admissibility regions of $M$.

A schematic constraint governing the field triple can be written as
\[
\partial_t \Phi + \nabla \cdot (\Phi\, \mathbf{v}) = -\lambda S,
\]
expressing the condition that admissibility density is transported by the trajectory
flow $\mathbf{v}$ and dissipated at a rate proportional to local entropy production.
Regions where entropy increases rapidly are regions where admissibility is consumed.
Stable function requires that $\Phi$ remains above a threshold $\Phi_{\min}$ even as
$S$ increases.

This equation can be derived as the Euler--Lagrange equation of an action functional
for the field triple. Define the RSVP Lagrangian density
\[
\mathcal{L}(\Phi, \mathbf{v}, S)
= \tfrac{1}{2}|\mathbf{v}|^2 - V(\Phi) - \alpha S + \beta\, \Phi\, \nabla \cdot \mathbf{v},
\]
where $V(\Phi)$ is a potential penalizing configurations below $\Phi_{\min}$,
$\alpha > 0$ weights the entropic cost of local randomization, and $\beta > 0$ couples
the admissibility field to the divergence of the trajectory flow. Variation with respect
to $\mathbf{v}$ yields the flow equation relating trajectory velocity to the gradient
of the admissibility potential; variation with respect to $\Phi$ yields the
conservation law stated above; and variation with respect to $S$ yields the condition
that entropy production is bounded by the coupling coefficient $\alpha$. The admissible
manifold $M$ emerges as the critical-point locus of this system under the constraint
that $S$ is non-decreasing along trajectories of $\mathbf{v}$. This variational
formulation makes the RSVP framework a genuine geometric field theory rather than a
merely interpretive framework: the admissible projection $\pi : X \to M$ is not
postulated but derived as the attractor of the constrained Euler--Lagrange flow.

\begin{definition}[RSVP Admissible Projection]
In the RSVP framework, the admissible projection $\pi : X \to M$ is the asymptotic map
\[
\pi(x) = \lim_{t \to \infty} \phi_t(x),
\]
where $\phi_t$ is the flow of $\mathbf{v}$ on $\mathcal{D}$, restricted to
$\{\Phi \geq \Phi_{\min}\}$ and conditioned on entropy non-decrease along trajectories.
The manifold $M$ is the attractor of this restricted flow.
\end{definition}

This recovers the admissible projection of Part~I as the long-time limit of constrained
flow. The trajectory $x$ is not preserved; what is preserved is the functional output
of the limiting point $\pi(x)$ on the admissible manifold.

\subsection{ARE as RSVP Field Dynamics}

In the ARE setting, the RSVP triple interprets as follows. The scalar field $\Phi$
encodes the probability density of ARE encodings over the group $G$. High-$\Phi$
regions correspond to encodings consistent with valid ARE transcripts: their aggregate
lies in the range of the decoder and produces a valid function output. The vector field
$\mathbf{v}$ is the flow induced by the aggregation operation. Individual encodings
$\hat{x}_i$ flow toward their aggregate $\hat{y} = \sum_i \hat{x}_i$. This flow is
constrained to preserve the functional invariant $f(x_1, \ldots, x_n)$ as a stable
attractor. The entropy field $S$ captures the randomness injected by local encoders.
High-$S$ regions of the encoding space are regions where many different joint input
configurations map to the same aggregate: these are the large-fiber regions of $\pi$,
corresponding to strong provenance destruction.

The security condition of ARE --- that the evaluator's view is simulatable from the
output alone --- is the RSVP condition that $S$ is high in the image of $\pi$. Many
microscopic configurations are consistent with any observed aggregate, and no efficient
procedure can distinguish them. Security and entropy are here not in tension. They are
the same phenomenon viewed from different angles.

\subsection{Bacterial Gene Expression as RSVP Field Dynamics}

In the bacterial gene expression setting, the RSVP triple interprets as follows. The
scalar field $\Phi$ encodes the spatial density of admissible configurations. High-$\Phi$
regions are regions of the cell where the relevant molecular interactions --- productive
transcription--translation coupling, RNase-transcript encounter --- can occur. The
nucleoid interior and the membrane periphery define distinct high-$\Phi$ regions for
different processes. The vector field $\mathbf{v}$ encodes the flow of molecular
trajectories: RNAP movement along DNA, ribosome translocation along mRNA, transcript
diffusion through the cytoplasm, RNase~E diffusion along the membrane. The entropy
field $S$ captures spatial and configurational randomness. High-$S$ regions are regions
of the cytoplasm where many different molecular trajectory histories produce the same
observable output. Low-$S$ regions are regions of tight coupling where the trajectory
is constrained to follow a specific path, as in the RNAP--ribosome coupled complex in
a high-TE gene.

The transition from coupled to uncoupled transcription, governed by translation
efficiency, is a transition in the RSVP sense: a change in the local structure of
$\mathbf{v}$ as the coupling flow becomes non-coherent, an increase in $S$ as the
spacer RNA becomes accessible to multiple degradation pathways, and a shift of the
trajectory from a low-entropy path on $M_{\mathrm{coupled}}$ to a higher-entropy path
on $M_{\mathrm{uncoupled}}$. The admissibility field $\Phi$ remains high in both
regimes, since functional regulation is achieved in either case, but the trajectory
by which the invariant is approached changes fundamentally.

\subsection{Obstruction, Cohomology, and the Topology of Non-Interaction}

The two central examples of admissible projection in this essay are geometrically dual
to each other in an important sense. In the ARE case, the projection destroys provenance
by mixing. The entropy injected by local encoders thermalizes the aggregate, making the
fiber of $\pi$ large. The functional invariant survives not because anything protects it
from noise, but because it is a stable functional of the aggregate that remains
well-defined regardless of which fiber element was realized.

In the bacterial case, the projection preserves function by separation. The spatial
admissibility constraint prevents certain molecular encounters from occurring, thereby
preventing certain degradation trajectories from being realized. The functional invariant
survives not because anything actively protects transcripts from RNase, but because the
geometry of the cell ensures that the relevant trajectories do not intersect.

These two mechanisms --- mixing and separation --- are dual in the sense of obstruction
cohomology. Mixing corresponds to a cohomological trivialization: the fiber of $\pi$ is
so large that no distinguisher can recover a global section from local data. The
extension problem --- recover $x$ from $\pi(x)$ --- has no solution because too many
solutions are consistent with the observed data. Separation corresponds to a
cohomological obstruction: the extension problem --- allow RNase to reach the nascent
transcript --- has no solution because the geometry prevents the required paths from
existing. The relevant cochains cannot be assembled into a global section.

Both mechanisms prevent the recovery of microscopic trajectory information. Both allow
functional invariants to be computed and preserved. They are two realizations of the
same abstract phenomenon: the impossibility of local-to-global extension in a
constrained topological setting.

\subsection{Life, Computation, and the RSVP Attractor}

The deepest unification offered by the RSVP framework is the identification of $M$
as an attractor of the dynamical system defined by the field triple. In both the
cryptographic and biological settings, the relevant systems are not designed to reach
$M$ in finite time and then halt. They are ongoing dynamical processes in which the
admissible manifold is continuously approached, partially realized, and partially
escaped. ARE encodings are not computed once; they are produced by parties who
participate in many future computations. Transcripts are not produced in isolation;
each transcript is one element of a continuous flux of production and degradation that
maintains cellular protein levels in a dynamical steady state.

In both cases, the attractor $M$ is not a fixed point but a statistical regime: a
region of the state space that the system spends most of its time near, and to which
it returns after perturbation. The functional invariants preserved by $\pi$ are not
static properties of a configuration; they are statistical properties of the ensemble
of configurations realized over time. A system computes, or a cell regulates, not by
instantiating a fixed microscopic configuration, but by maintaining a dynamical
ensemble of microscopic trajectories whose time-averaged projection $\pi(x(t))$ remains
concentrated on the admissible manifold $M$. Function is not a property of any single
trajectory. It is a property of the projection of a trajectory ensemble onto an
admissible invariant manifold under constrained entropy flow.

%------------------------------------------------------
\section{Part V: Provenance, Privacy, and the Topology of Accessibility}
%------------------------------------------------------

\subsection{Privacy as Accessibility Geometry}

Classical notions of privacy are framed negatively. A system is private if an observer
cannot infer a hidden variable, cannot distinguish between two neighboring databases,
or cannot reconstruct a secret from observable outputs. These formulations treat privacy
as an epistemic limitation imposed upon an observer. The admissible projection framework
developed in this essay suggests a different interpretation. Privacy is not fundamentally
about secrecy. It is about accessibility geometry.

In the ARE setting, the evaluator fails to reconstruct the microscopic inputs not because
the inputs are hidden behind an encrypted wall, but because the admissible projection
$\pi : X \to M$ collapses many microscopic configurations into the same observable
invariant. The evaluator is trapped on the quotient manifold. The inaccessible
information is inaccessible because the topology of the fibers forbids unique
reconstruction.

In the bacterial setting, the same structure appears physically. RNase~E cannot access
the nascent transcript because the trajectories required for encounter do not exist
within the admissible geometry of the cell. The transcript is not encrypted. It is
spatially inaccessible. No key is needed to protect it because no path leads to it.

\begin{definition}[Topological Privacy]
A system exhibits \emph{topological privacy} with respect to an observable $f : X \to
Y$ if the admissible projection $\pi : X \to M$ induces fibers $\pi^{-1}(m)$ whose
internal structure is inaccessible to an observer constrained to trajectories within
$M$. The inaccessibility need not be cryptographic; it may arise because the admissible
geometry forbids the trajectories required for reconstruction.
\end{definition}

Under this interpretation, privacy becomes a property of the topology of accessibility
rather than merely a property of information encoding. The abstract structure appears
identically in shuffle-based computation, in membrane-localized degradation, in
thermodynamic coarse-graining, and in cognitive abstraction. In each case the observer
or agent is restricted to a low-dimensional projection manifold whose geometry forbids
reconstruction of the microscopic trajectory ensemble.

\subsection{Fiber Structure and Functional Equivalence Classes}

The fibers of $\pi$ play a central structural role throughout this essay. They are the
equivalence classes of microscopic trajectories that produce the same admissible
invariant. In the ARE setting, the fiber over a transcript $\hat{y}$ is
\[
\pi^{-1}(\hat{y})
=
\bigl\{
(x_1,\ldots,x_n,r_1,\ldots,r_n)
\;\big|\;
\textstyle\sum_i \mathrm{Enc}(x_i;r_i)=\hat{y}
\bigr\}.
\]
This fiber is enormous. Many distinct joint inputs and randomness assignments produce
the same aggregate transcript. Security follows from the inability to distinguish which
fiber element was realized.

In the biological setting, a measured protein abundance level may correspond to many
distinct transcription initiation histories, many distinct ribosome loading sequences,
many distinct degradation trajectories, and many distinct diffusion paths. The
observable survives because it is fiber-invariant: its value is the same regardless
of which microscopic realization within the fiber was actually traversed.

\begin{proposition}[Invariant Stability Principle]
The stability of a functional observable increases as its dependence on
fiber-internal structure decreases. An observable constant across large fibers of $\pi$
remains stable under substantial microscopic variation.
\end{proposition}

This principle is immediate. Any observable depending sensitively on microscopic
trajectory detail will fluctuate strongly under perturbation of the fiber structure.
The principle explains simultaneously why shuffle MPC remains secure under random
permutation, why biological regulation tolerates molecular stochasticity, why
thermodynamic observables remain stable under microscopic chaos, and why semantic
cognition tolerates representational variation. The observable survives because it is a
property of the equivalence class, not of any individual trajectory.

\subsection{The Algebra of Admissible Observables}

The fiber structure of $\pi$ determines not merely which observables are stable but
which observables are \emph{definable} on the quotient manifold $M$. This observation
motivates a precise characterization of the admissible observable algebra.

\begin{definition}[Admissible Observable Algebra]
Let $\pi : X \to M$ be an admissible projection. The \emph{algebra of admissible
observables} is the set
\[
\mathcal{O}(M) = \bigl\{ f : X \to \mathbb{R} \;\big|\; f(x) = f(x') \text{ whenever }
\pi(x) = \pi(x') \bigr\}.
\]
Elements of $\mathcal{O}(M)$ are precisely the functions on $X$ that factor through
$\pi$: they are constant on fibers and therefore well-defined on the quotient.
\end{definition}

This definition has immediate structural content. The algebra $\mathcal{O}(M)$ is
closed under pointwise addition, multiplication, and composition with continuous
functions, making it a genuine commutative algebra. Any measurement performed on a
system governed by the projection $\pi$ that produces a stable, reproducible result
must belong to $\mathcal{O}(M)$; measurements depending on fiber-internal structure
cannot be stably reproduced across repeated experiments that realize different fiber
elements.

The analogy with gauge theory is precise here. In a gauge theory, physical observables
are those that are invariant under gauge transformations --- the group acting on the
fibers of the principal bundle. The algebra of gauge-invariant observables is the
appropriate algebra of physical quantities. In the present framework, the equivalence
relation $x \sim x'$ whenever $\pi(x) = \pi(x')$ plays the role of the gauge group,
and $\mathcal{O}(M)$ plays the role of the gauge-invariant observable algebra. The
admissible projection $\pi$ is the quotient map of this gauge structure.

In the ARE setting, $\mathcal{O}(M)$ consists precisely of those functions of the
aggregate transcript $\hat{y}$ that can be computed without access to the individual
encodings. Simulation-based security is exactly the statement that any efficient
computation on the evaluator's view belongs to $\mathcal{O}(M)$: the only efficiently
computable function of the aggregate is the designated output $f(x_1, \ldots, x_n)$.

In the bacterial setting, $\mathcal{O}(M)$ consists of those cellular observables that
are independent of the exact microscopic trajectory ensemble: protein abundance, mRNA
half-life, the coupling coefficient between transcription and translation. These are the
quantities that remain stable under molecular stochasticity, spatial heterogeneity, and
thermodynamic fluctuation. Quantities not in $\mathcal{O}(M)$ --- such as the exact
identity of which ribosome translates which transcript, or the precise spatial path of
a given mRNA molecule --- are not stably accessible to any measurement on the
biological system operating at steady state.

The classical epistemological claim that the map is not the territory becomes, under
this framework, a structural statement with precise content: the territory is $X$, the
map is $\pi$, and the only quantities accessible to a system operating on the admissible
manifold are those belonging to $\mathcal{O}(M)$. What the system \emph{can} know is
determined not by the richness of $X$ but by the structure of $\pi$. This is not
merely epistemic incompleteness. The system itself only preserves the quotient geometry.
The map is dynamically primary because the dynamical attractor \emph{is} the projection
class, not any element of its preimage.

\subsection{The Symmetry Between Mixing and Separation}

The two instances of topological privacy examined in this essay are related by a
structural symmetry. Mixing enlarges fibers by aggregating many inputs into the same
output. Separation eliminates trajectories by making certain regions of the state space
geometrically unreachable. In both cases the observer is denied access to the
microscopic interior. In mixing, the interior is crowded with indistinguishable
alternatives. In separation, the interior is empty of reachable paths.

This symmetry suggests that topological privacy is not a single mechanism but a class
of mechanisms unified by the common property that the fiber structure of the admissible
projection is inaccessible to trajectories confined to the image. The specific
realization --- whether by entropy injection, spatial separation, temporal gating, or
structural obstruction --- is secondary to the abstract topological property it
instantiates.

In the RSVP language, mixing corresponds to a large entropy field $S$ in the image of
$\pi$, while separation corresponds to a steep gradient in the admissibility field
$\Phi$ that falls to zero across the boundary between admissible and inadmissible
regions. Both configurations produce topological privacy. They are dual in the sense
that each can be expressed as a limiting case of the other under appropriate
deformations of the field triple $(\Phi, \mathbf{v}, S)$.

%------------------------------------------------------
\section{Part VI: Constraint, Entropy, and Persistent Computation}
%------------------------------------------------------

\subsection{Constraint Before Content}

The ARE and bacterial systems examined in this essay both exhibit a striking inversion
of classical informational ontology. Traditionally, one imagines that content exists
first, that computation manipulates that content, and that constraints are externally
imposed afterward as a kind of friction. The systems studied here exhibit the opposite
structure. The constraints determine which trajectories can exist before any stable
functional content appears. The admissibility geometry precedes the invariant.

In the ARE setting, the evaluator cannot reconstruct provenance because the shuffle and
aggregation structure eliminate the relevant degrees of freedom before decoding occurs.
The admissible geometry of the protocol determines what kinds of observables are even
definable. No transcript produced by valid ARE encoders can carry individual input
information; the structure of the encoding scheme makes such a transcript
topologically impossible. Content --- in the sense of function output --- is the
residue left after the constraint geometry has eliminated everything else.

In the bacterial setting, RNase accessibility, transcriptional coupling, and degradation
kinetics are governed by spatial topology before any particular transcript fate is
realized. The geometry determines which molecular trajectories can become stable. A
transcript encodes the same genetic sequence regardless of whether it is coupled to a
ribosome or prematurely released into the cytoplasm, but its fate --- and therefore the
functional content it can contribute to the cell --- is determined entirely by geometry.

\begin{remark}
This inversion is central to the RSVP framework. The stable object is not the
microscopic configuration itself, but the persistence class induced by the admissible
projection. Objects emerge from constrained flow. They are not fundamental primitives
that later encounter constraints; they are the residue of constraint satisfaction.
\end{remark}

\subsection{Persistent Computation Without Persistent Representation}

One of the deepest implications of the ARE result is that universal computation does
not require persistent local representation. The classical symbolic model of computation
assumes stable memory registers, identifiable state transitions, recoverable
intermediate representations, and explicit communication pathways. The shuffle/ARE
framework undermines all four assumptions simultaneously.

The evaluator does not know which message came from which party. It cannot reconstruct
intermediate representations. It observes only an aggregated projection. And yet it
still computes arbitrary functions. This is not merely secure computation. It is
computation surviving representational destruction.

The same phenomenon appears biologically. The cell does not preserve exact transcript
trajectories, exact diffusion histories, exact ribosome-polymerase alignments, or exact
degradation pathways. Nevertheless, functional regulation persists robustly across
thousands of simultaneous transcription and translation events, across the stochastic
variation inherent in molecular encounter dynamics, and across the spatial heterogeneity
of the cellular interior.

\begin{definition}[Persistent Computation]
A system performs \emph{persistent computation} if it maintains stable functional
invariants despite continual destruction, randomization, or inaccessibility of
microscopic representational structure. Persistent computation is computation whose
stability arises precisely because the invariant does not depend on microscopic
recoverability.
\end{definition}

Persistent computation is therefore not computation with noise added afterward. It is
computation organized from the outset around invariants that are indifferent to the
microscopic paths by which they are reached. The admissible projection $\pi$ is not an
approximation; it is the exact structure of the computation.

\subsection{Entropy as Functional Selector}

Classical thermodynamic intuition often treats entropy as antagonistic to order and
functionality. The systems studied here demonstrate a more nuanced picture. In ARE
systems, increased entropy strengthens privacy and stabilizes the functional invariant
by enlarging the fibers of $\pi$. Entropy destroys provenance while preserving
computation. More randomness, injected locally, means larger fibers, which means
stronger security and more robust invariant preservation. Entropy is here a resource,
not a threat.

In weakly coupled bacterial systems, excessive entropy destabilizes the coupled
transcriptional manifold by allowing trajectory divergence, Rho loading, and premature
termination. The entropy that flows into the spacer between RNAP and ribosome is
entropy that disrupts coupling, collapses the coupled trajectory, and sends the
transcript into a high-degradation regime. Entropy here is a threat, not a resource.

Entropy therefore has no universal sign. Its effect on function depends entirely on
the geometry of admissibility relative to the direction of entropic flow. The
appropriate formulation is
\[
\text{Persistence} \neq \text{low entropy}.
\]
Rather,
\[
\text{Persistence} = \text{stability of invariant structure under entropy flow}.
\]
A system remains functional not by minimizing entropy universally, but by ensuring that
entropy flow remains tangent to the admissible manifold $M$. When entropy flow exits the
tangent bundle of $M$, the invariant destabilizes. In the ARE setting, the injected
entropy remains within the security-preserving quotient geometry; it enlarges fibers
without perturbing the invariant. In the bacterial setting, low translation efficiency
permits entropy flow into trajectory sectors that are transverse to the coupled
manifold, thereby destabilizing the coupling that maintains transcript integrity.

The RSVP framework accommodates both behaviors naturally. The entropy field $S$ is not
intrinsically destructive or constructive. Its effect depends on the interaction between
the scalar admissibility field $\Phi$, the trajectory flow field $\mathbf{v}$, and the
topology of the projection manifold $M$. When $\nabla S$ points into the tangent bundle
of $M$, entropy increase is compatible with invariant stability. When $\nabla S$ points
transversely to $M$, entropy increase destabilizes the projection.

%------------------------------------------------------
\section{Part VII: Toward a General Theory of Admissible Systems}
%------------------------------------------------------

\subsection{Beyond Two Domains}

The examples examined in this essay suggest that admissible projection is not a
specialized mechanism restricted to cryptography or bacterial regulation. It may instead
be a general organizational principle for complex systems in which functional stability
must be maintained under conditions of partial irreversibility and continuous microscopic
variation.

In morphogenesis, embryological development proceeds through massive irreversible
projection. Cellular trajectories branch, signaling pathways diverge, epigenetic states
become inaccessible, and yet stable anatomical invariants emerge with extraordinary
reliability. The mature organism is not a preservation of microscopic developmental
history. It is a stable projection class produced by constrained developmental flow
under conditions of continuous entropy production. The full trajectory ensemble of every
cell from zygote to adult is irretrievably lost; the organism is the invariant that
survives this loss.

In cognition, perceptual systems discard vast quantities of sensory detail. Semantic
systems collapse many distinct representations into invariant categories. Memory systems
preserve compressed attractors rather than exact trajectories. The invariant that the
cognitive system computes --- the recognition of an object, the parsing of a sentence,
the application of a concept --- does not depend on exact reconstruction of the sensory
or representational path. It is a property of the equivalence class, accessible through
many distinct microscopic trajectories. The mind computes by projection, and its
reliability arises from the same source as the reliability of ARE systems: the
functional invariant is constant across large fibers, and therefore robust to variation
within them.

\subsection{The Generalized Principle of Functional Persistence}

The central principle emerging from this essay can now be stated in its general form.

\begin{proposition}[Generalized Functional Persistence]
Let $X$ be a microscopic trajectory space and let $\pi : X \to M$ be an admissible
projection preserving a class of observables $\mathcal{F}$. A system exhibits
persistent function if the observables in $\mathcal{F}$ are constant on sufficiently
large fibers of $\pi$; if entropy flow increases microscopic indistinguishability
without destabilizing the admissible manifold $M$; and if perturbations of microscopic
trajectories remain confined to equivalence classes preserving $\mathcal{F}$.
\end{proposition}

This principle unifies privacy-preserving computation, bacterial gene regulation,
thermodynamic coarse-graining, semantic abstraction, and morphogenesis under a single
geometric description. The common structure is not symbolic manipulation. It is
admissible invariant persistence under irreversible projection. The principle does not
depend on the nature of the microscopic substrate; it holds whenever an admissible
projection exists with the stated properties.

The generalization also clarifies what distinguishes robust systems from fragile ones.
A fragile system is one whose functional output depends on fiber-internal structure:
small changes in the microscopic trajectory produce large changes in the observable.
A robust system is one whose functional output is genuinely fiber-invariant: the same
output is produced by a large and topologically connected class of microscopic
realizations. Robustness, under this formulation, is not a property added on top of
function. It is a consequence of the geometric relationship between the functional
observable and the fiber structure of the admissible projection.

\subsection{Admissible Projection and Renormalization}

The structure of admissible projection is closely related to the logic of renormalization
in theoretical physics, and making this connection explicit substantially broadens the
scope of the framework. In the Wilsonian approach to renormalization, a quantum or
statistical field theory is not described by a single action functional but by a family
of effective theories parametrized by a cutoff scale $\Lambda$. Integrating out degrees
of freedom below the scale $\Lambda$ produces a coarser effective theory whose
predictions for long-wavelength observables match those of the original theory,
while the short-wavelength microscopic structure becomes inaccessible.

The admissible projection $\pi : X \to M$ can be interpreted as a generalized
renormalization operator. Introduce a scale parameter $\Lambda$ and define a
scale-dependent family of projections
\[
\pi_\Lambda : X \longrightarrow M_\Lambda,
\]
where $M_\Lambda$ is the admissible manifold at scale $\Lambda$: the quotient of $X$
by the equivalence relation that identifies all microscopic configurations
indistinguishable by any observable of characteristic scale larger than $\Lambda$.
As $\Lambda$ decreases, $M_\Lambda$ becomes coarser: fewer microscopic distinctions
survive, and the fibers of $\pi_\Lambda$ grow. As $\Lambda$ increases, $M_\Lambda$
becomes finer: more microscopic structure is retained, and the projection approaches
the identity on $X$.

Under this formulation, persistent observables at scale $\Lambda$ are precisely the
elements of $\mathcal{O}(M_\Lambda)$: functions that are constant on the fibers of
$\pi_\Lambda$. The flow from fine to coarse scale --- decreasing $\Lambda$ --- is a
flow in the space of projections, and the renormalization group equation governing this
flow is structurally analogous to the RSVP admissibility conservation law. Renormalization
fixed points correspond to admissible manifolds $M_\Lambda$ that are self-similar
under further coarsening: projections whose fiber structure does not change under
small decreases in $\Lambda$. These are the most robust observables, stable not merely
under a fixed projection but under the entire renormalization flow.

In the ARE setting, the relevant scale is the information scale: how much
distinguishing power the evaluator has access to. At any finite information scale, the
evaluator's accessible observables form a subalgebra of $\mathcal{O}(M)$. Security
corresponds to the condition that the designated output $f(x_1, \ldots, x_n)$ is the
only element of $\mathcal{O}(M)$ accessible at the evaluator's scale. In the bacterial
setting, the relevant scales are spatial and temporal: cellular observables at the
scale of protein abundance and mRNA lifetime are renormalization fixed structures of
the underlying molecular dynamics, stable precisely because the fiber structure of
the projection from molecular trajectories to cellular observables is preserved under
the biological analogue of scale flow.

\subsection{Morphisms of Admissible Systems}

The admissible projection framework admits a natural categorical formulation that
makes precise what it means for two systems to be structurally equivalent, and what
it means for one system to simulate another.

Define the category $\mathbf{Adm}$ whose objects are triples $(X, \pi, M)$
consisting of a microscopic trajectory space $X$, an admissible manifold $M$, and
an admissible projection $\pi : X \to M$. A morphism
\[
F : (X, \pi_X, M_X) \longrightarrow (Y, \pi_Y, M_Y)
\]
in $\mathbf{Adm}$ is a pair of maps $(F_X : X \to Y,\; \tilde{F} : M_X \to M_Y)$
such that the diagram
\[
\begin{array}{ccc}
X & \xrightarrow{F_X} & Y \\
\downarrow \pi_X & & \downarrow \pi_Y \\
M_X & \xrightarrow{\tilde{F}} & M_Y
\end{array}
\]
commutes: $\pi_Y \circ F_X = \tilde{F} \circ \pi_X$. The commutativity condition
requires that $F_X$ maps fiber elements of $\pi_X$ to fiber elements of $\pi_Y$:
if two microscopic configurations are equivalent under $\pi_X$, their images under
$F_X$ are equivalent under $\pi_Y$. Morphisms in $\mathbf{Adm}$ therefore preserve
the admissible observable algebra: if $f \in \mathcal{O}(M_Y)$ then
$f \circ \tilde{F} \in \mathcal{O}(M_X)$, meaning that any observation definable
on the target system pulls back to a well-defined observation on the source system.

This categorical formulation makes precise several structural relationships that
the essay has treated informally. The ARE scheme is a morphism from the trajectory
space of joint inputs and randomness $(X_{\mathrm{ARE}}, \pi_{\mathrm{ARE}},
M_{\mathrm{ARE}})$ to the output space, with the encoding map playing the role of
$F_X$ and the decoder playing the role of $\tilde{F}$. Security is the condition
that this morphism has no nontrivial sections: there is no map from $M_{\mathrm{ARE}}$
back to $X_{\mathrm{ARE}}$ compatible with the categorical structure.

The biological regulatory system is a morphism from the space of molecular trajectory
ensembles to the space of cellular observables. The membrane-localization constraint
on RNase~E is a condition on the morphism $F_X$: it restricts the image of $F_X$ to
trajectory ensembles in which degradation encounters are kinetically inaccessible,
thereby ensuring that the induced map $\tilde{F}$ on the admissible manifold sends
coupled transcriptional states to stable protein-production states.

Isomorphisms in $\mathbf{Adm}$ --- morphisms with invertible $\tilde{F}$ --- correspond
to systems that preserve the same admissible invariant algebra through different
microscopic mechanisms. The three bacterial species examined in Part~III are
isomorphic in $\mathbf{Adm}$ with respect to the functional observable of controlled
gene expression, even though their microscopic implementations differ entirely.
This categorical isomorphism is what the essay has called species-specific admissibility
topology: different objects in $\mathbf{Adm}$ with the same admissible invariant structure.

\subsection{Failure Modes: Catastrophic Projection Collapse}

The essay has so far presented admissible projection primarily in its successful form:
systems in which functional invariants survive microscopic collapse. An equally
important part of the framework is the characterization of failure. Not all irreversible
projections are admissible. When projection becomes too coarse, functional invariants
are destroyed rather than preserved. This is what we call catastrophic projection
collapse.

Formally, catastrophic collapse occurs when the projection $\pi$ fails to respect the
functional structure of the system. The precise condition is the existence of
$x, x' \in X$ such that
\[
f(x) \neq f(x') \quad \text{but} \quad \pi(x) = \pi(x').
\]
That is, two microscopically distinct configurations that produce different functional
outputs are mapped to the same point on the admissible manifold. The projection
conflates functionally distinct states. This is not the benign fiber-enlargement
of secure systems; it is a destructive coarsening that makes the invariant
unrecoverable from the projected state.

In the ARE setting, catastrophic collapse corresponds to a decoding failure: two
different function values produce indistinguishable aggregates. This is precisely
what correctness requirements forbid. The ARE correctness condition
\[
\Pr\bigl[\mathrm{Dec}(\hat{y}) = f(x_1, \ldots, x_n)\bigr] \geq 1 - \varepsilon
\]
is the quantitative bound on how rarely catastrophic collapse is permitted to occur.
Perfect correctness would require $\varepsilon = 0$, eliminating catastrophic
collapse entirely; statistical correctness with negligible $\varepsilon$ bounds
its probability.

In the biological setting, catastrophic collapse corresponds to the failure of
cellular regulation. When the admissibility constraint breaks down --- when
$\Gamma_{\mathrm{RNase}} \cap \Gamma_{\mathrm{nascent}}$ is no longer approximately
empty, as in organisms with cytoplasmic RNase~E or in experimental relocalization
mutants --- co-transcriptional degradation increases and the functional invariant
of transcript stability becomes unreliable. The organism does not cease to function
immediately, because the admissibility field $\Phi$ falls below $\Phi_{\min}$ only
locally and transiently, but the robustness of the system against perturbation is
reduced. A sufficiently severe disruption of the admissibility geometry would
produce catastrophic collapse of the coupled transcriptional manifold.

The distinction between productive irreversibility and catastrophic collapse is therefore
the central criterion for evaluating any admissible system. Productive irreversibility
enlarges fibers while preserving fiber invariance: the projection $\pi$ maps more
microscopic configurations to the same admissible output, but those configurations
are functionally equivalent. Catastrophic collapse conflates functionally distinct
configurations: the projection fails to separate $f$-preimages. The design of robust
systems consists precisely in ensuring that entropy flow remains within the former
regime and does not cross into the latter.

The deepest philosophical implication of the admissible projection framework is that
microscopic ontology loses explanatory primacy. If stable function depends primarily on
equivalence classes, admissibility geometry, and invariant-preserving projection, then
the exact microscopic realization becomes secondary not merely epistemically but
structurally.

This does not imply that microscopic structure is unreal. It implies that persistent
systems are organized around projection classes rather than exact trajectories, and that
the explanatorily relevant level of description is the level of the admissible manifold
rather than the level of the microscopic configuration space. The ARE evaluator does not
need the original messages; the bacterial cell does not preserve transcript histories;
the cognitive system does not preserve raw sensory trajectories; the thermodynamic
system does not preserve microscopic particle paths. What persists is the invariant.

Classical scientific ontology tends to privilege the microscopic as the locus of
ultimate explanation. The tradition from Democritus through Boltzmann to molecular
biology treats the macroscopic as supervening on the microscopic in a way that makes
the macroscopic derivative and the microscopic foundational. The admissible projection
framework inverts this priority not by denying the existence of microscopic structure
but by demonstrating that it is the projection --- the map $\pi : X \to M$ --- and not
the preimage that is explanatorily fundamental for any system whose outputs are
persistent and computable.

The stability that matters for function is not the stability of trajectories but the
stability of projection classes. The invariant that survives collapse is not located in
any particular microscopic realization. It is located in the geometry of the map itself.
Stability belongs not to the microscopic realization but to the projection class.

This observation can be consolidated as a general structural principle.

\begin{proposition}[Projection Robustness Principle]
A dynamical system is robust precisely to the extent that its operational observables
factor through an admissible projection $\pi : X \to M$ whose fibers absorb microscopic
perturbations without altering invariant structure. That is, a system is robust with
respect to a class of perturbations $\delta x$ if and only if for all $x \in X$ and
all admissible perturbations $\delta x$, $\pi(x + \delta x) = \pi(x)$: perturbations
remain within the fiber and do not change the projection class.
\end{proposition}

This proposition unifies the robustness properties of the two systems examined in this
essay. Shuffle MPC is robust under any permutation of messages, under any choice of
randomness assignments by the local encoders, and under any substitution within the
fiber $\pi^{-1}(\hat{y})$, because all such variations preserve the aggregate $\hat{y}$
and therefore leave $\pi(x)$ unchanged. Bacterial gene regulation is robust under
molecular stochasticity, diffusion noise, and individual transcript fate variation,
because these perturbations do not move the system off the admissible manifold $M$:
the statistical observables, protein levels, mRNA half-lives, and coupling ratios,
remain fiber-invariant under the relevant perturbation classes. Fragility enters
precisely when a perturbation crosses a fiber boundary, moving the system to a
different projection class and thereby altering the invariant structure. In the
bacterial setting this corresponds to the transition from coupled to uncoupled
transcription; in the ARE setting it would correspond to a break in the encoding
structure that allows the evaluator to distinguish fiber elements.

\begin{remark}
This is the sense in which function survives collapse. The collapse is not a failure of
the system. It is the mechanism by which persistent invariants become possible. A system
that preserved all microscopic detail would be a system in which every perturbation
threatens every output. A system organized around admissible projection tolerates
perturbation precisely because the invariant does not depend on the perturbed element.
The collapse creates the robustness.
\end{remark}

%------------------------------------------------------
\section{Part VIII: Seed Development as Admissible Signalling Geometry}
%------------------------------------------------------

\subsection{The Developing Seed as a System of Partially Isolated Dynamical Domains}

A recent single-nucleus RNA sequencing atlas of early \textit{Arabidopsis thaliana} seed
development \citep{martin2026atlas} provides a third major empirical instance of the
admissible projection framework, one that extends its reach from computation and
prokaryotic gene regulation into multicellular developmental biology. The study profiles
54,210 nuclei across 3, 5, and 7 days after pollination, capturing the transcriptional
programs of the embryo, endosperm, and seed coat at high resolution. What the atlas
reveals is not a centrally orchestrated developmental program but a distributed system
of partially isolated signalling domains maintaining developmental coherence across
constrained interfaces.

The three main seed tissues are genetically distinct and symplastically separated:
they cannot exchange arbitrary intracellular state through direct cytoplasmic continuity.
Coordination therefore does not arise from unrestricted informational transparency
but from sparse signalling across well-defined interfaces. This is the structural
signature of an admissible system. The full microscopic biochemical state of each
tissue is inaccessible to the others; what propagates across tissue boundaries is a
low-dimensional projection sufficient to maintain global developmental invariants.

In the language of Part~I, the developing seed instantiates an admissible projection
\[
\pi : X_{\mathrm{seed}} \longrightarrow M_{\mathrm{seed}},
\]
where $X_{\mathrm{seed}}$ is the full microscopic state space of all transcriptional,
hormonal, and signalling dynamics across all tissues, and $M_{\mathrm{seed}}$ is the
admissible developmental manifold: the space of tissue coordination states consistent
with successful seed maturation. The critical observation is that
$M_{\mathrm{seed}} \neq \bigcup_i M_i$ for the individual tissue manifolds $M_i$.
The global developmental outcome is not the union of isolated tissue programs; it
emerges from their constrained interaction across admissibility interfaces.

\subsection{Short Secreted Peptides as Low-Dimensional Coordination Operators}

The atlas identifies the endosperm, particularly its micropylar and chalazal regions,
as the primary hub for expression of short secreted peptides (SSPs), including
members of the DEFL, LCR, and LTP families, as well as specifically characterized
peptides including TWS1, ALE1, KRS, and RALFL3 \citep{martin2026atlas}. Many of
these are upregulated after fertilization and are concentrated at the critical tissue
interfaces where the symplastically isolated compartments must coordinate.

The SSPs are not carrying exhaustive developmental information across tissue boundaries.
They are carrying constraint updates: low-dimensional signals sufficient to update
the admissible developmental trajectory of the receiving tissue without requiring
full informational transparency into the sending tissue's internal state.

This is the developmental analogue of the ARE encoding. A thermostat does not
communicate the full microscopic thermodynamic state of a building. It communicates
a single low-dimensional control signal. Likewise, the SSPs appear to coordinate
admissibility boundaries --- thresholds of developmental readiness, timing of
transitions, compatibility of nutrient-flow states --- rather than transmitting full
biochemical descriptions. The receiving tissue updates its trajectory on $M_i$ in
response to a low-dimensional signal from the interface, without accessing the
internal state of the neighboring compartment.

This can be formalized as a coordination operator
\[
\mathcal{C}_{ij} : M_i \longrightarrow M_j,
\]
mapping the projected developmental state of tissue $i$ to an update of the admissible
state of tissue $j$. In the category $\mathbf{Adm}$ introduced in Part~VII, these
coordination operators are morphisms between the individual tissue admissible systems.
The SSPs instantiate these morphisms physically: they are the molecular realization
of low-bandwidth inter-manifold communication. Their sequence conservation across
species, and the evidence for rapid evolution at tissue interfaces under conditions of
genetic conflict, is precisely what one would predict of morphisms under selection
pressure to maintain inter-manifold coordination while allowing microscopic variation
within each tissue's internal state.

\subsection{The Chalazal Endosperm as a Developmental Interface Manifold}

The chalazal endosperm (CZE) is the region of the endosperm closest to the vascular
supply of the maternal tissue. The atlas reveals that the CZE is not a homogeneous
compartment. It exhibits apical-basal transcriptional polarity within its coenocytic
cyst: the gene AT3G49307 marks the apical region while RALFL3 marks the basal founder
nuclei, which show transcriptional similarity to the adjacent maternal tissues,
the chalazal proliferating tissue and the inner integument \citep{martin2026atlas}.
Pseudotime analysis suggests the basal cyst follows an independent developmental
trajectory from the rest of the CZE.

This result is of particular significance for the admissible projection framework
because it demonstrates that even within a nominally single tissue compartment,
signalling geometry remains anisotropic and directionally structured. The CZE is not
merely a tissue. It is a developmental interface manifold, a geometric region where
multiple admissibility flows converge, interact under directional constraints, and
separate toward distinct developmental attractors.

In RSVP terms, the chalazal endosperm instantiates a non-trivial vector field
$\mathbf{v}(x, t)$ within the tissue interior. Developmental trajectories are not
isotropic diffusion but directed constrained transport along geometry-dependent
admissibility gradients. The apical-basal polarity is a manifestation of this
directionality: the entropy field $S$ and the admissibility density $\Phi$ have
non-zero gradients within the cyst, producing distinct trajectory bundles in its
apical and basal regions.

The basal founder nuclei, with their transcriptional similarity to adjacent maternal
tissues, are especially interpretable through this lens. They occupy a region of the
developmental manifold where the fiber structure of the projection $\pi$ is shared with
neighboring tissue compartments: the same admissible developmental states are
realizable from more than one microscopic tissue identity. This is fiber overlap across
tissue boundaries, and it is precisely what enables the chalazal region to function as
an interface: it mediates between the endosperm's internal developmental trajectory
and the maternal tissue's signalling geometry through a shared region of
$M_{\mathrm{seed}}$.

\subsection{Brassinosteroid Localization and Curvature of Developmental Admissibility Space}

One of the most spatially precise results in the atlas is the localization of
brassinosteroid biosynthesis. The genes BR6OX1 and BR6OX2, encoding the final enzymes
in the brassinosteroid synthesis pathway, are expressed with high specificity in a
micropylar subtype of the outer integument \citep{martin2026atlas}. Brassinosteroid
signalling, as reflected by expression of the BZR1/BES1 family transcription factors,
is then detected in the micropylar endosperm at 5 days after pollination.

The spatial restriction of hormone production to a single specialized cell subtype
represents a concrete biological realization of a non-uniform admissibility field.
Brassinosteroid is not distributed homogeneously across the seed. It is produced
at a localized source and establishes a spatial gradient in developmental admissibility
space. In RSVP terms,
\[
\nabla \Phi \neq 0
\]
across the seed interior: the admissibility density $\Phi$ has a spatial structure
determined by the geometry of the hormone source, not by centralized global instruction.
Developmental behavior in neighboring tissues depends on position relative to this
gradient rather than on access to any symbolic specification of what to do.

This is the developmental analogue of a gravitational or electromagnetic field. The
brassinosteroid-producing micropylar outer integument behaves as a localized field
generator whose influence propagates outward through concentration gradients, inducing
tissue-specific responses according to the local value of $\Phi$ rather than through
explicit messaging. The receiving tissues in the micropylar endosperm do not need to
know the internal state of the outer integument; they respond to the projected
admissibility signal encoded in the local hormone concentration.

\subsection{Development as Constraint Propagation Across Partially Isolated Manifolds}

The atlas as a whole supports a reinterpretation of plant embryogenesis that replaces
the implicit genetic blueprint metaphor with a constraint-propagation picture. The
blueprint model implies that developmental outcomes are encoded in a centrally
accessible program executed sequentially. The atlas evidence supports a different
picture: global developmental coherence emerges from local constraint updates
propagated across partially isolated tissue domains through sparse signalling interfaces.

No single tissue appears to contain or execute the complete developmental program for
the seed. The embryo, endosperm, and seed coat each operate according to their own
internal trajectories on their respective admissibility manifolds $M_i$. Global
coordination emerges not from any tissue accessing the internal state of the others,
but from the low-dimensional coordination signals --- the SSPs, the hormonal gradients,
the mechanical and nutrient flow cues --- that update admissibility boundaries across
interfaces.

The analogy to flocking behavior is exact. No bird contains the blueprint for the
flock. Global structure emerges from local interaction constraints that each bird
can compute from its immediate neighbourhood. Likewise, each seed tissue responds
to locally accessible signals while the global seed developmental manifold emerges
from their constrained interaction. The developmental system solves a distributed
coordination problem without centralizing information, precisely as the ARE system
solves a secure computation problem without centralizing state.

The rapid evolution observed at maternal-filial tissue interfaces, concentrated in
extracellular domains and intrinsically disordered regions of secreted proteins
\citep{martin2026atlas}, is the evolutionary fingerprint of this architecture.
Interfaces under genetic conflict --- where maternal and paternal interests in resource
allocation diverge --- are precisely the locations where the coordination operators
$\mathcal{C}_{ij}$ are under the strongest selection pressure to change. The molecular
machinery implementing inter-manifold communication evolves rapidly; the global
developmental invariant it maintains evolves much more slowly. Function survives at
the level of $M_{\mathrm{seed}}$ while representation diverges at the level of the
individual morphisms. This is the evolutionary expression of the Projection Robustness
Principle: perturbations to the molecular implementation of coordination operators
remain within fibers of the developmental projection, leaving the invariant structure
of $M_{\mathrm{seed}}$ intact.

%------------------------------------------------------
\section{Part IX: The Formal Isomorphism Between Shuffle Privacy and Spatial Regulation}
%------------------------------------------------------

\subsection{From Structural Analogy to Categorical Equivalence}

The preceding parts of this essay have established that shuffle-model secure computation
and spatial gene regulation in bacteria share the same abstract structure: both are
instances of admissible projection $\pi : X \to M$ in which fiber-constant observables
persist while microscopic provenance is destroyed. However, the claim of structural
unity remains incomplete so long as the two systems are treated as analogous rather
than formally equivalent. The purpose of this part is to close that gap by constructing
a rigorous isomorphism, within the category $\mathbf{Adm}$, between the shuffle
operator governing ARE and the spatial accessibility operator governing RNase~E--transcript
interaction.

The strategy proceeds in five steps. First, the symmetry group of each system is
identified and shown to act on the respective fiber bundle in structurally identical
ways. Second, the projection operator for each system is characterized as a specific
type of constraint Hamiltonian, and the two Hamiltonians are shown to satisfy the same
structural equation. Third, the spatial collision integral formalizing the biological
projection is derived and placed in correspondence with the aggregation functional of
the ARE. Fourth, the obstruction cohomology class preventing admissible liftability is
computed for both systems and shown to lie in the same cohomological position. Fifth,
the invariant subalgebra lemma and the dynamical attractor theorem are proved in the
general setting, completing the isomorphism.

\subsection{Symmetry Groups and Fiber Bundle Structure}

The first step in constructing the isomorphism is to identify the symmetry group of
each system. In the ARE shuffle, the relevant symmetry group is the symmetric group
$S_n$ acting on the set of labeled inputs. Let $X_{\mathrm{ARE}}$ be the space of
all $n$-tuples $(x_1, \ldots, x_n, r_1, \ldots, r_n)$ of inputs and randomness.
The shuffle operator $\sigma$ is a map $\sigma : X_{\mathrm{ARE}} \to M_{\mathrm{ARE}}$
where $M_{\mathrm{ARE}} = G$ is the Abelian group over which encodings are summed.
The group $S_n$ acts on $X_{\mathrm{ARE}}$ by permuting the labels:
\[
\rho \cdot (x_1, \ldots, x_n, r_1, \ldots, r_n)
= (x_{\rho(1)}, \ldots, x_{\rho(n)}, r_{\rho(1)}, \ldots, r_{\rho(n)}),
\quad \rho \in S_n.
\]
The aggregate sum $\hat{y} = \sum_i \mathrm{Enc}(x_i; r_i)$ is manifestly
permutation-invariant: $\sigma(\rho \cdot x) = \sigma(x)$ for all $\rho \in S_n$.
The shuffle operator is therefore a quotient map for the $S_n$-action. The fiber
$\sigma^{-1}(\hat{y})$ over any aggregate $\hat{y} \in G$ is the $S_n$-orbit of any
element of $X_{\mathrm{ARE}}$ that maps to $\hat{y}$, together with all randomness
assignments producing the same sum.

In the bacterial gene expression setting, the relevant symmetry group is the isometry
group $\mathrm{Iso}(3) = \mathbb{R}^3 \rtimes O(3)$ of Euclidean space, acting on the
space $X_{\mathrm{bio}}$ of molecular trajectory ensembles. Let $\tau : X_{\mathrm{bio}}
\to M_{\mathrm{bio}}$ be the spatial accessibility projection, where $M_{\mathrm{bio}}
= \{0, 1\}$ encodes the binary outcome of transcript stability. The membrane-targeting
sequence of RNase~E constrains the enzyme to a region within distance $d_{\min}$ of the
inner membrane. Any isometry $g \in \mathrm{Iso}(3)$ that preserves this distance
constraint --- that is, any $g$ mapping the membrane-proximal region to itself ---
also preserves the accessibility partition of $X_{\mathrm{bio}}$:
\[
\tau(g \cdot \gamma) = \tau(\gamma)
\quad \text{for all } \gamma \in X_{\mathrm{bio}},\; g \in \mathrm{Stab}(\Omega_{\mathrm{RNase}}),
\]
where $\Omega_{\mathrm{RNase}} \subset \mathbb{R}^3$ is the membrane-proximal region
accessible to RNase~E and $\mathrm{Stab}(\Omega_{\mathrm{RNase}})$ is the stabilizer
subgroup of $\mathrm{Iso}(3)$ preserving this region. The spatial projection $\tau$
is thus a quotient map for the action of this stabilizer.

\begin{definition}[Structural Symmetry of an Admissible System]
An admissible system $(X, \pi, M)$ has \emph{structural symmetry group} $G_\pi$ if
$G_\pi$ acts on $X$ by a group action $\cdot : G_\pi \times X \to X$ such that
$\pi(g \cdot x) = \pi(x)$ for all $g \in G_\pi$ and all $x \in X$. The fibers of $\pi$
are then exactly the orbits of $G_\pi$.
\end{definition}

The ARE system has structural symmetry group $G_\sigma = S_n$. The bacterial system has
structural symmetry group $G_\tau = \mathrm{Stab}(\Omega_{\mathrm{RNase}})
\leq \mathrm{Iso}(3)$. These groups are not isomorphic as abstract groups, but they
play the same structural role: each acts freely on the fibers of the respective
projection, making the fiber over any base point $m \in M$ a torsor for the symmetry
group. It is this structural role, not the group-theoretic identity, that defines the
isomorphism.

\subsection{Constraint Hamiltonians and the Projection Mechanism}

The second step is to show that both projections arise from the same type of constraint
Hamiltonian. A constraint Hamiltonian is a function $H_{\mathrm{cons}} : X \to \mathbb{R}$
whose level sets define the admissibility partition. The projection $\pi$ maps each
$x \in X$ to the level set $\{x' : H_{\mathrm{cons}}(x') = H_{\mathrm{cons}}(x)\}$
quotiented by the symmetry group action.

For the ARE system, the constraint Hamiltonian is the additive aggregate:
\[
H_{\mathrm{cons}}^{\mathrm{ARE}}(x_1, \ldots, x_n, r_1, \ldots, r_n)
= \sum_{i=1}^n \mathrm{Enc}(x_i; r_i) \in G.
\]
The level sets of $H_{\mathrm{cons}}^{\mathrm{ARE}}$ are exactly the fibers of $\sigma$.
The "collapse" of the ARE system consists in the evaluator accessing only the value of
$H_{\mathrm{cons}}^{\mathrm{ARE}}$, not the individual terms in the sum. Since
$H_{\mathrm{cons}}^{\mathrm{ARE}}$ is invariant under $S_n$, its level sets are unions
of $S_n$-orbits, and the admissible observable algebra $\mathcal{O}(M_{\mathrm{ARE}})$
consists exactly of functions of $H_{\mathrm{cons}}^{\mathrm{ARE}}$.

For the bacterial system, the constraint Hamiltonian is the encounter potential:
\[
H_{\mathrm{cons}}^{\mathrm{bio}}(\gamma)
= \int_0^T \mathbf{1}[\gamma(t) \in \Omega_{\mathrm{RNase}}]\, dt,
\]
where $\gamma : [0, T] \to \mathbb{R}^3$ is a trajectory of the nascent transcript
and $\mathbf{1}[\gamma(t) \in \Omega_{\mathrm{RNase}}]$ is the indicator that the
transcript is within the RNase-accessible region at time $t$. When
$H_{\mathrm{cons}}^{\mathrm{bio}}(\gamma) = 0$, the transcript has zero encounter
probability and projects to the stable outcome; when
$H_{\mathrm{cons}}^{\mathrm{bio}}(\gamma) > \theta$ for a threshold $\theta$,
the transcript projects to the degraded outcome. The level sets of
$H_{\mathrm{cons}}^{\mathrm{bio}}$ are the fibers of $\tau$, invariant under the
stabilizer of $\Omega_{\mathrm{RNase}}$.

Both Hamiltonians have the form of a \emph{constraint accumulator}: they sum or
integrate a local quantity (the encoding contribution, or the accessibility indicator)
over the relevant index set (parties, or time), and the admissible projection is the
level-set map of this accumulator. The structural equation satisfied by both is:
\[
\pi(x) = \pi(x') \iff H_{\mathrm{cons}}(x) = H_{\mathrm{cons}}(x'),
\]
which is the defining condition for $\pi$ to be the quotient map of the constraint
Hamiltonian's level-set partition.

\subsection{The Spatial Collision Integral}

The biological constraint Hamiltonian $H_{\mathrm{cons}}^{\mathrm{bio}}$ can be made
more precise through the theory of diffusion-limited reactions. The encounter probability
between a nascent transcript at position $\mathbf{r}_0$ and RNase~E distributed
according to a density $\rho_{\mathrm{RNase}}(\mathbf{r})$ on the inner membrane is
given by the path-integral expression
\[
P(\text{encounter} \mid \gamma)
= 1 - \exp\!\left(
-\kappa \int_0^T \rho_{\mathrm{RNase}}(\gamma(t))\, dt
\right),
\]
where $\kappa$ is the intrinsic reaction rate and $\gamma(t)$ is the transcript's
spatial trajectory. The expected encounter probability over the ensemble of transcript
trajectories is
\[
\bar{P}(\text{encounter})
= 1 - \mathbb{E}_\gamma\!\left[\exp\!\left(
-\kappa \int_0^T \rho_{\mathrm{RNase}}(\gamma(t))\, dt
\right)\right].
\]
When RNase~E is membrane-localized, $\rho_{\mathrm{RNase}}(\mathbf{r})$ is concentrated
on the membrane surface $\partial\Omega_{\mathrm{cell}}$. For a transcript whose
trajectory $\gamma$ is confined to the nucleoid interior
$\Omega_{\mathrm{nucleoid}} \subset \Omega_{\mathrm{cell}}$ at distance greater than
$d_{\min}$ from the membrane, $\rho_{\mathrm{RNase}}(\gamma(t)) \approx 0$ for all $t$,
giving $\bar{P}(\text{encounter}) \approx 0$. This is the spatial admissibility
constraint in quantitative form.

The spatial projection $\tau : X_{\mathrm{bio}} \to M_{\mathrm{bio}}$ is then defined
by
\[
\tau(\gamma) =
\begin{cases}
0 & \text{(stable)} \quad \text{if } \bar{P}(\text{encounter} \mid \gamma) < \varepsilon, \\
1 & \text{(degraded)} \quad \text{if } \bar{P}(\text{encounter} \mid \gamma) \geq \varepsilon,
\end{cases}
\]
for a threshold $\varepsilon$ reflecting the minimum encounter probability required
for degradation to proceed at a biologically significant rate. The fiber
$\tau^{-1}(0)$ consists of all transcript trajectories for which the time-averaged
density of RNase~E along the path is below the degradation threshold.

\begin{proposition}[Structural Equivalence of Constraint Accumulators]
Let $H_{\mathrm{cons}}^{\mathrm{ARE}}$ and $H_{\mathrm{cons}}^{\mathrm{bio}}$ be as
defined above. Both satisfy the structural equation
\[
\pi(x) = \pi(x') \iff H_{\mathrm{cons}}(x) = H_{\mathrm{cons}}(x'),
\]
and both are invariant under their respective structural symmetry groups. The admissible
projection in each case is the level-set quotient map of the constraint accumulator.
\end{proposition}

\begin{proof}
For the ARE system: $H_{\mathrm{cons}}^{\mathrm{ARE}}$ is a group homomorphism from
$X_{\mathrm{ARE}}$ (under componentwise encoding) to $G$. Its level sets are the fibers
of $\sigma$. Invariance under $S_n$ follows from the commutativity of $G$: permuting
the summands does not change the sum. The quotient map is $\sigma$ by construction.

For the biological system: $H_{\mathrm{cons}}^{\mathrm{bio}}(\gamma) = \kappa \int_0^T
\rho_{\mathrm{RNase}}(\gamma(t))\,dt$ is a functional of the trajectory $\gamma$. Its
level sets define $\tau$. Invariance under $\mathrm{Stab}(\Omega_{\mathrm{RNase}})$
follows from the fact that any isometry preserving $\Omega_{\mathrm{RNase}}$ also
preserves $\rho_{\mathrm{RNase}}$ pointwise, hence preserves the integral. The level-set
map is $\tau$ by definition.

In both cases the projection is surjective onto $M$ (every output value is achievable),
the fibers are non-trivial (multiple elements of $X$ map to each base point), and the
symmetry group acts freely on fibers, making each fiber a torsor. The structural
equations are identical in form.
\end{proof}

\subsection{The Invariant Subalgebra Lemma}

Having established that both systems are governed by the same structural equation for
their constraint accumulators, we now prove the invariant subalgebra lemma: any
observable not invariant under the system's structural symmetry group is uncomputable
by an observer restricted to $M$.

\begin{proposition}[Invariant Subalgebra Lemma]
Let $(X, \pi, M)$ be an admissible system with structural symmetry group $G_\pi$.
An observable $f : X \to \mathbb{R}$ is computable by an observer restricted to $M$
if and only if $f$ is $G_\pi$-invariant, that is, $f(g \cdot x) = f(x)$ for all
$g \in G_\pi$ and all $x \in X$.
\end{proposition}

\begin{proof}
($\Rightarrow$) Suppose $f$ is computable from $M$, meaning there exists
$\tilde{f} : M \to \mathbb{R}$ with $f = \tilde{f} \circ \pi$. For any $g \in G_\pi$
and $x \in X$:
\[
f(g \cdot x) = \tilde{f}(\pi(g \cdot x)) = \tilde{f}(\pi(x)) = f(x),
\]
where the second equality uses $G_\pi$-invariance of $\pi$. Hence $f$ is $G_\pi$-invariant.

($\Leftarrow$) Suppose $f$ is $G_\pi$-invariant. We must show $f$ descends to $M$.
Define $\tilde{f}(m) = f(x)$ for any $x \in \pi^{-1}(m)$. This is well-defined: if
$x, x' \in \pi^{-1}(m)$ then $\pi(x) = \pi(x') = m$, so $x' = g \cdot x$ for some
$g \in G_\pi$ (since the fibers are $G_\pi$-orbits), and $G_\pi$-invariance of $f$
gives $f(x') = f(g \cdot x) = f(x)$. Thus $\tilde{f}$ is a well-defined function on $M$
satisfying $f = \tilde{f} \circ \pi$.
\end{proof}

\begin{remark}
For the ARE system, $G_\pi = S_n$ and the invariant observables are exactly the
symmetric functions of the encodings, which by the ARE security condition reduce to
functions of the aggregate sum alone. For the bacterial system,
$G_\pi = \mathrm{Stab}(\Omega_{\mathrm{RNase}})$ and the invariant observables are
exactly those depending only on the time-averaged encounter probability --- that is,
on whether the transcript is stable or degraded, not on the specific trajectory by
which it reached that state.
\end{remark}

\subsection{Obstruction to Admissible Liftability}

We now prove the central structural theorem: neither system admits an admissible
section, that is, a map $s : M \to X$ compatible with $\pi$ that would allow an
observer on $M$ to recover the microscopic trajectory.

\begin{definition}[Admissible Section]
A \emph{admissible section} of $\pi : X \to M$ is a map $s : M \to X$ such that
$\pi \circ s = \mathrm{id}_M$ and $s$ is compatible with the admissible observable
algebra: for every $f \in \mathcal{O}(M)$, $f \circ \pi \circ s = f$.
\end{definition}

\begin{proposition}[Obstruction to Liftability]
Let $(X, \pi, M)$ be an admissible system with structural symmetry group $G_\pi$
acting freely and non-trivially on the fibers of $\pi$. Then no admissible section
$s : M \to X$ exists.
\end{proposition}

\begin{proof}
Suppose for contradiction that an admissible section $s$ exists. Fix any $m \in M$ and
let $x_0 = s(m) \in \pi^{-1}(m)$. Since $G_\pi$ acts freely and non-trivially on the
fiber $\pi^{-1}(m)$, there exists $g \in G_\pi$, $g \neq e$, such that $g \cdot x_0
\neq x_0$. But $\pi(g \cdot x_0) = \pi(x_0) = m$, so both $x_0$ and $g \cdot x_0$ are
elements of $\pi^{-1}(m)$.

Now consider the function $f : X \to \mathbb{R}$ defined by $f(x) = d(x, x_0)$ for some
metric $d$ on $X$. This function distinguishes $x_0$ from $g \cdot x_0$: $f(x_0) = 0$
while $f(g \cdot x_0) = d(g \cdot x_0, x_0) > 0$. However, $f$ is not $G_\pi$-invariant
(since $f(g \cdot x_0) \neq f(x_0)$), so by the invariant subalgebra lemma, $f \notin
\mathcal{O}(M)$ and is uncomputable from $M$.

Any admissible section $s$ would select a specific element $s(m) = x_0$ from the fiber
$\pi^{-1}(m)$. But the selection criterion cannot be expressed in terms of observables
from $\mathcal{O}(M)$ alone, since distinguishing $x_0$ from $g \cdot x_0$ requires a
function outside $\mathcal{O}(M)$. Therefore no admissible selection rule exists, and
$s$ cannot be defined in a way compatible with the admissible observable algebra.

This is the obstruction. It is precisely the cohomological statement that the fibration
$\pi : X \to M$ has no admissible section: the first Čech cohomology class
$[\pi] \in \check{H}^1(M, \underline{G_\pi})$ associated with the principal
$G_\pi$-bundle is non-trivial, preventing the existence of a global section.
\end{proof}

\begin{remark}
For the ARE system, the non-triviality of $[\sigma] \in \check{H}^1(M_{\mathrm{ARE}},
\underline{S_n})$ is the formal statement that the shuffle cannot be inverted: knowing
the aggregate $\hat{y}$ does not determine which of the $n!$ possible permutations of
inputs produced it, and no admissible procedure can select a canonical representative
of the fiber. For the bacterial system, the non-triviality of $[\tau] \in
\check{H}^1(M_{\mathrm{bio}}, \underline{G_\tau})$ is the formal statement that
knowing a transcript is stable does not determine which of the many admissible
nucleoid-interior trajectories it followed.
\end{remark}

The central conclusion follows immediately.

\begin{proposition}[Privacy and Regulation as Failure of Admissible Liftability]
\[
\text{privacy/regulation} = \text{failure of admissible liftability}.
\]
More precisely: the security of the ARE shuffle and the stability of the bacterial
regulatory system are both equivalent to the non-existence of an admissible section
of their respective projection maps, which is in turn equivalent to the non-triviality
of the associated Čech cohomology class.
\end{proposition}

\subsection{The Isomorphism Theorem}

We can now state and prove the main theorem of this part: the ARE shuffle system and
the bacterial spatial accessibility system are isomorphic objects in $\mathbf{Adm}$
with respect to their functional invariant structure, even though their structural
symmetry groups are not group-theoretically isomorphic.

\begin{definition}[Functional Isomorphism in $\mathbf{Adm}$]
Two admissible systems $(X_1, \pi_1, M_1)$ and $(X_2, \pi_2, M_2)$ are
\emph{functionally isomorphic} if there exists an isomorphism $\tilde{F} : M_1 \to M_2$
in $\mathbf{Adm}$ such that the induced map on admissible observable algebras
$\tilde{F}^* : \mathcal{O}(M_2) \to \mathcal{O}(M_1)$ is an algebra isomorphism,
and both systems have non-trivial Čech cohomology obstruction classes preventing
admissible liftability.
\end{definition}

\begin{proposition}[Functional Isomorphism of ARE and Bacterial Systems]
The ARE admissible system $(X_{\mathrm{ARE}}, \sigma, M_{\mathrm{ARE}})$ and the
bacterial spatial admissible system $(X_{\mathrm{bio}}, \tau, M_{\mathrm{bio}})$
are functionally isomorphic in $\mathbf{Adm}$ with respect to the following invariant
structure: both have a designated computable functional observable (the function output
$f(x_1, \ldots, x_n)$ for ARE; the protein production rate for biology) that is the
unique element of their respective admissible observable algebras $\mathcal{O}(M)$
that an observer on $M$ can stably compute, and both have non-trivial obstruction
classes preventing recovery of microscopic trajectory information.
\end{proposition}

\begin{proof}
The proof proceeds by exhibiting the isomorphism $\tilde{F}$ explicitly and verifying
its properties.

Define $\tilde{F} : M_{\mathrm{ARE}} \to M_{\mathrm{bio}}$ by mapping the subset of
$M_{\mathrm{ARE}}$ corresponding to valid function outputs to the stable state
$0 \in M_{\mathrm{bio}}$, and mapping all other aggregate values to the unstable state.
More precisely, for the binary output case where $f(x_1, \ldots, x_n) \in \{0,1\}$:
the aggregate $\hat{y}$ with $\mathrm{Dec}(\hat{y}) = 0$ maps to $\tilde{F}(\hat{y}) =
\text{stable}$, and $\hat{y}$ with $\mathrm{Dec}(\hat{y}) = 1$ maps to
$\tilde{F}(\hat{y}) = \text{degraded}$.

Under this map, the designated computable observable in ARE --- the function output
$f(x_1, \ldots, x_n)$ --- corresponds to the designated computable observable in the
biological system --- the binary stability state. The admissible observable algebras
$\mathcal{O}(M_{\mathrm{ARE}})$ and $\mathcal{O}(M_{\mathrm{bio}})$ are both one-dimensional
in the relevant subspace (one designated output observable, one binary state), and
$\tilde{F}^*$ sends the stability observable to the function output observable, establishing
the algebra isomorphism on the designated observables.

Both systems have non-trivial Čech obstruction classes by the Obstruction to Liftability
proposition: the ARE shuffle has class $[\sigma] \in \check{H}^1(M_{\mathrm{ARE}},
\underline{S_n})$ and the bacterial system has class $[\tau] \in \check{H}^1(M_{\mathrm{bio}},
\underline{G_\tau})$, both non-trivial by the freeness and non-triviality of the
respective group actions on fibers.

The functional isomorphism is therefore established: both systems preserve exactly the
same abstract structure (one designated computable observable, non-trivial liftability
obstruction, fiber-invariant algebra), through different physical mechanisms (permutation
symmetry and additive aggregation for ARE; spatial symmetry and diffusion geometry for
biology) that play the same structural role in $\mathbf{Adm}$.
\end{proof}

\subsection{The Dynamical Attractor Theorem}

The final component of the formal isomorphism is the convergence theorem: under the
RSVP field dynamics, all trajectories in $X$ satisfying the admissibility constraints
converge to $M$. This gives the isomorphism its dynamical rather than merely static
content.

\begin{proposition}[Dynamical Attractor Theorem]
Let $(\Phi, \mathbf{v}, S)$ be the RSVP field triple on a domain $\mathcal{D}$
with Lagrangian density $\mathcal{L} = \tfrac{1}{2}|\mathbf{v}|^2 - V(\Phi) -
\alpha S + \beta\, \Phi\, \nabla \cdot \mathbf{v}$. Suppose the constraint potential
$V(\Phi)$ is strictly convex on $\{\Phi \geq \Phi_{\min}\}$ with a unique minimum on
the admissible manifold $M$. Then under the Euler--Lagrange flow of $\mathcal{L}$
with entropy non-decrease, every trajectory $\phi_t(x)$ with $\Phi(\phi_t(x)) \geq
\Phi_{\min}$ converges to $M$ as $t \to \infty$.
\end{proposition}

\begin{proof}
The Euler--Lagrange equation for $\mathbf{v}$ from $\mathcal{L}$ gives
\[
\frac{D\mathbf{v}}{Dt} = -\nabla V(\Phi) + \beta\,\nabla(\Phi\, \nabla \cdot \mathbf{v}),
\]
where $D/Dt$ is the material derivative along the flow. The admissibility conservation
equation $\partial_t \Phi + \nabla \cdot (\Phi\, \mathbf{v}) = -\lambda S$ implies
that $\Phi$ decreases along trajectories proportionally to local entropy production.

Define the Lyapunov function $\mathcal{V}(x) = V(\Phi(x)) + \alpha S(x)$. Along
any trajectory $\phi_t(x)$:
\[
\frac{d}{dt}\mathcal{V}(\phi_t(x))
= \nabla V(\Phi) \cdot \partial_t \Phi + \alpha\, \partial_t S.
\]
From the admissibility conservation equation, $\partial_t \Phi = -\nabla \cdot
(\Phi\,\mathbf{v}) - \lambda S \leq -\lambda S$ (since $\nabla \cdot (\Phi\,\mathbf{v})
\geq 0$ in the admissible regime by the flow structure). The entropy non-decrease
condition gives $\partial_t S \geq 0$.

For $\Phi$ near but above $\Phi_{\min}$, strict convexity of $V$ on the admissible
region implies $\nabla V(\Phi) \cdot (-\lambda S) \leq 0$ with equality only on $M$
(where $V$ achieves its minimum). Combined with the entropy non-decrease, $\mathcal{V}$
is a decreasing function along trajectories except on $M$. By the LaSalle invariance
principle applied to the Euler--Lagrange flow, all trajectories in the admissible
region $\{\Phi \geq \Phi_{\min}\}$ converge to the largest invariant subset of
$\{\frac{d}{dt}\mathcal{V} = 0\}$, which by strict convexity is exactly $M$.
\end{proof}

\begin{remark}
This theorem applies to both the ARE and bacterial systems under their respective RSVP
interpretations. For ARE, the Lyapunov function $\mathcal{V}$ measures the deviation
of the encoding distribution from the security-preserving quotient geometry; the
Euler--Lagrange flow drives the encoding distribution toward the maximal-entropy fiber
structure on $M_{\mathrm{ARE}}$. For the bacterial system, $\mathcal{V}$ measures the
deviation of the transcript trajectory ensemble from the membrane-separated admissible
region; the RSVP flow drives transcript trajectories toward the stable nucleoid-interior
attractor basin on $M_{\mathrm{bio}}$.
\end{remark}

\subsection{Synthesis: The Unified Structural Picture}

The results of this part can now be assembled into a single unified structural picture.
Both the ARE shuffle system and the bacterial spatial regulation system are admissible
systems in $\mathbf{Adm}$ with the following common properties, established by the
preceding propositions and theorems.

Their projections arise from constraint Hamiltonians of the same structural type: they
are both constraint accumulators, additive over an index set, invariant under a
structural symmetry group acting freely on fibers. Their admissible observable algebras
consist exactly of the invariant subalgebras of their respective symmetry groups.
No admissible section exists for either: both carry non-trivial Čech cohomology
obstruction classes that prevent microscopic recovery from manifold-level observation.
Both are dynamical attractors under RSVP field evolution: the Lyapunov function
$\mathcal{V}$ decreases to zero on $M$ under the constrained Euler--Lagrange flow.
They are functionally isomorphic in $\mathbf{Adm}$: there exists an isomorphism of
their designated observable structures that respects the admissible algebra and the
cohomological obstruction.

The formal conclusion is that privacy in shuffle-model computation and spatial
regulation in bacterial gene expression are not merely analogous but are isomorphic
instances of the same abstract phenomenon: admissible projection with non-trivial
liftability obstruction, preserving a designated computable invariant while making
microscopic provenance unrecoverable. The physical mechanisms differ entirely, but the
structural position of those mechanisms within the category $\mathbf{Adm}$ is identical.

%------------------------------------------------------
\section{Interlude: The Three-Layer Architecture and Bubble Boundaries}
%------------------------------------------------------

\subsection{Separating the Foundational, Dynamical, and Realization Layers}

The material developed in Parts~I through~IX operates simultaneously across three
distinct levels of abstraction, and the coherence of the monograph as a whole benefits
from making these levels explicit before proceeding. The three layers are: the
foundational geometric layer, the dynamical layer, and the domain realization layer.
Each layer has its own vocabulary, its own formal objects, and its own mode of
argument; the power of the framework derives from the fact that all three are
structurally consistent instantiations of the same underlying geometry.

The \emph{foundational geometric layer} consists of the objects and relations
that require no dynamical or empirical specification to be defined. Its primary
components are the admissible projection $\pi : X \to M$, the admissible invariant
manifold $M$, the admissible observable algebra $\mathcal{O}(M)$, the structural
symmetry group $G_\pi$, the obstruction class $[\pi] \in \check{H}^1(M, \underline{G_\pi})$,
and the bubble boundary $\partial B$ defined below. These objects are axiomatic in the
sense that they can be specified independently of any particular physical or
computational realization.

The \emph{dynamical layer} consists of the RSVP field triple $(\Phi, \mathbf{v}, S)$
and the associated Lagrangian $\mathcal{L} = \tfrac{1}{2}|\mathbf{v}|^2 - V(\Phi) -
\alpha S + \beta\,\Phi\,\nabla\cdot\mathbf{v}$, its Euler--Lagrange flow, and the
collapse operators that govern how systems move through admissibility space.
This layer explains the temporal evolution of admissible systems: it answers the
question of how a system approaches $M$, how it remains on $M$ under perturbation,
and how it exits $M$ under catastrophic collapse. The RSVP framework is the dynamical
layer. The Spherepop calculus governs discrete collapse events and irreversible
quotient formation within this layer; the Yarncrawler framework governs the
transport geometry connecting distinct admissibility basins.

The \emph{domain realization layer} consists of specific instantiations of the
foundational and dynamical structures in empirical or computational systems.
Each domain realization names concrete objects in the abstract framework: for ARE,
$X$ is the input-randomness space, $\pi$ is the shuffle operator, and $G_\pi = S_n$;
for bacterial gene regulation, $X$ is the molecular trajectory ensemble, $\pi$ is the
spatial accessibility projection, and $G_\pi = \mathrm{Stab}(\Omega_{\mathrm{RNase}})$;
for seed development, $X$ is the space of tissue-level biochemical states,
$\pi$ is the inter-tissue coordination projection, and the coordination operators
$\mathcal{C}_{ij}$ are the morphisms in $\mathbf{Adm}$. Additional domain realizations
will be developed in Parts~X and beyond, including typography, semantic cognition,
and generative computation.

The relationship between layers is: the foundational layer provides the static
categorical scaffolding; the dynamical layer provides the flow equations governing
how systems traverse that scaffolding; and the realization layer interprets specific
physical or computational systems within it. A theorem proved at the foundational layer
holds in all domain realizations. A result derived at the dynamical layer applies
wherever the RSVP field equations govern the relevant flows. A claim made at the
realization layer holds for that specific domain and may or may not generalize.

\subsection{Bubble Boundaries and the Geometry of Collapse}

The most important geometric object not yet given an explicit definition in the main
text is the bubble boundary $\partial B$. This object formalizes the intuition that
admissible systems occupy a region of trajectory space with a definite boundary,
crossing which constitutes collapse.

\begin{definition}[Admissibility Bubble]
Let $(X, \pi, M)$ be an admissible system with structural symmetry group $G_\pi$.
The \emph{admissibility bubble} $B \subseteq X$ is the maximal open set of
configurations that project to admissible outputs under small perturbations:
\[
B = \bigl\{ x \in X \;\big|\; \exists\,\varepsilon > 0 : \pi(x + \delta x) = \pi(x)
\text{ for all } \|\delta x\| < \varepsilon \bigr\}.
\]
The \emph{bubble boundary} is
\[
\partial B = \bigl\{ x \in X \;\big|\; \forall\,\varepsilon > 0,\,
\exists\,\delta x \text{ with } \|\delta x\| < \varepsilon
\text{ and } \pi(x + \delta x) \neq \pi(x) \bigr\}.
\]
\end{definition}

Elements of $B$ are robust under infinitesimal perturbation: small changes in the
microscopic trajectory do not alter the projection class. Elements of $\partial B$ are
marginally stable: arbitrarily small perturbations can change the projection class.
Elements of $X \setminus (B \cup \partial B)$ are already outside the admissible
regime.

The bubble boundary is the geometric locus of incipient collapse. In the Projection
Robustness Principle proved in Part~VII, the condition $\pi(x + \delta x) = \pi(x)$
for all admissible perturbations is precisely the condition that $x \in B$. Fragility
is the condition $x \in \partial B$. Catastrophic projection collapse, as defined in
the failure modes subsection, is the transition $x \to x + \delta x$ with
$x + \delta x \notin B \cup \partial B$.

\begin{proposition}[Bubble Boundary as Entropy Tangency Threshold]
Under the RSVP field dynamics, $x \in \partial B$ if and only if $\nabla S(x)$ acquires
a non-zero component transverse to the tangent bundle of $M$ at $\pi(x)$:
\[
x \in \partial B \iff (\nabla S(x))_\perp \neq 0,
\]
where $(\nabla S)_\perp$ denotes the component of $\nabla S$ in the normal bundle
$N(M) = TX|_M / TM$.
\end{proposition}

\begin{proof}
By the RSVP Lagrangian, the admissibility density $\Phi$ decreases at rate
$\lambda S$ along trajectories. The condition $x \in B$ requires that this
decrease is tangential to $M$: the projection class $\pi(x)$ remains constant under
the flow. This tangentiality condition is $\partial_t \Phi|_M = -\lambda S|_M$
remaining in $TM$, which holds if and only if $\nabla S$ has no transverse component
to $M$. When $(\nabla S)_\perp \neq 0$, the entropy gradient drives the admissibility
density off $M$, causing the projection class to change, placing $x$ at or beyond
$\partial B$.
\end{proof}

This proposition is the precise form of the entropy tangency condition mentioned
throughout the essay. Entropy is destructive not in itself but in its direction:
entropy flow tangent to $M$ enlarges fibers without altering projection classes;
entropy flow transverse to $M$ drives the system toward $\partial B$ and into collapse.

The bubble boundary formalism subsumes many phenomena discussed in earlier parts.
Premature transcription termination in bacteria is the crossing of $\partial B$ by the
RNAP--ribosome coupled trajectory, driven by transverse entropy flow from the spacer
RNA into the Rho-loading regime. Privacy leakage in ARE is the crossing of $\partial B$
by an encoding distribution, allowing an adversary to distinguish fiber elements.
Developmental failure in the seed system is the collapse of coordination operators
$\mathcal{C}_{ij}$ when tissue interface geometry places the signalling system at
$\partial B$. Each of these failure modes is a boundary-crossing event in the same
geometric sense.

\subsection{The Four-Layer Stack: RSVP, Spherepop, Yarncrawler, and Adm}

The full theoretical architecture of this monograph involves four interacting frameworks
that operate at complementary levels of description. Making their relationships precise
is necessary for the monograph's coherence as it expands into new domains.

RSVP provides the \emph{continuous field dynamics}: the equations governing how the
triple $(\Phi, \mathbf{v}, S)$ evolves over time and space, how admissibility density
is transported and dissipated, and how the constrained Euler--Lagrange flow drives
trajectories toward attractor basins on $M$. RSVP is the physics of the dynamical layer.

Spherepop provides the \emph{collapse topology}: the discrete event structure governing
irreversible quotient formation, bubble stabilization, and the moment of projection
collapse. Where RSVP tracks the continuous approach to $\partial B$, Spherepop governs
what happens at and after the crossing: the Pop operator collapses a bubble into its
projection class, the Bind operator establishes a new equivalence relation, the Collapse
operator finalizes the quotient, and the Refuse operator prevents certain transitions.
Spherepop is the topology of irreversible events within the dynamical layer.

Yarncrawler provides the \emph{inter-bubble transport geometry}: the mechanisms by
which information, constraint updates, or coordination signals are transmitted between
distinct admissibility basins. Where RSVP governs motion within a single admissibility
regime and Spherepop governs the transition between regimes, Yarncrawler governs
the network structure connecting multiple regimes. In the seed development context,
Yarncrawler corresponds to the signalling channels carrying SSP coordination operators
between tissue compartments. In the cognitive context, it corresponds to the
associative pathways connecting distinct semantic attractor basins.

The category $\mathbf{Adm}$ provides the \emph{categorical abstraction layer}: the
language in which all of the above can be stated domain-independently. Objects in
$\mathbf{Adm}$ are tuples $(X, \pi, M, \mathcal{O})$; morphisms are commutative
diagrams preserving observable algebras; functors between $\mathbf{Adm}$ and other
categories formalize the relationships between different domain realizations; natural
transformations between such functors capture the structural equivalences that the
isomorphism theorems of Part~IX make precise.

The four layers are not competing accounts but complementary ones: RSVP provides the
differential geometry, Spherepop provides the combinatorial topology, Yarncrawler
provides the graph-theoretic transport structure, and $\mathbf{Adm}$ provides the
categorical abstraction that lets results proved in any one domain transfer to all
structurally isomorphic ones.

%------------------------------------------------------
\section{Part X: Typography as Admissible Projection Geometry}
%------------------------------------------------------

\subsection{Glyphs as Invariant Manifolds}

The theory of admissible projection, developed in abstract form in Parts~I through~IX,
finds one of its most immediately visualizable instantiations in typography. A glyph
is not a bitmap. It is not a Bézier curve specification. It is not a set of stroke
parameters. A glyph is an equivalence class: the set of all possible geometric
realizations that an observer would recognize as the same symbolic identity. The
letter A, in any font, in any weight, in any orientation permitting recognition, in
any rendering medium from ink to pixel to chalk dust, is the same element of a
semantic projection manifold. The glyph is the projection class, not any of its
realizations.

This observation is not merely a metaphor for the abstract framework. It is a direct
instance of it. Define the glyph trajectory space $X_{\mathrm{stroke}}$ to be the
high-dimensional space of all possible geometric realizations of a single glyph:
all possible stroke trajectories, curvature profiles, weight distributions, spatial
frequency components, and rendering parameters. Define the semantic glyph manifold
$M_{\mathrm{semantic}}$ to be the low-dimensional space of recognizable symbolic
identities. The glyph projection is
\[
\pi_{\mathrm{glyph}} : X_{\mathrm{stroke}} \longrightarrow M_{\mathrm{semantic}},
\]
mapping each geometric realization to its equivalence class under perceptual recognition.
The fiber $\pi_{\mathrm{glyph}}^{-1}(A)$ over the symbol $A$ consists of all geometric
realizations that an observer projects onto the symbolic identity $A$: Roman,
bold, italic, condensed, extended, handwritten, graffitied, engraved, rendered in
fog, approximated in ASCII --- these are all elements of the same fiber.

The admissible observable algebra $\mathcal{O}(M_{\mathrm{semantic}})$ consists of
exactly those functions of glyph geometry that depend only on the projection class:
functions that are constant on fibers of $\pi_{\mathrm{glyph}}$. Semantic identity
is in $\mathcal{O}(M_{\mathrm{semantic}})$; exact stroke curvature at a specific point
is not. This is precisely the distinction that makes reading possible: human perceptual
systems compute elements of $\mathcal{O}(M_{\mathrm{semantic}})$ without accessing
fiber-internal detail.

\begin{definition}[Glyph Admissibility Bubble]
The \emph{admissibility bubble} $B_g \subset X_{\mathrm{stroke}}$ of a glyph
$g \in M_{\mathrm{semantic}}$ is the maximal open set of geometric realizations
$x \in X_{\mathrm{stroke}}$ such that
\[
\pi_{\mathrm{glyph}}(x + \delta x) = g
\quad \text{for all sufficiently small perturbations } \delta x.
\]
The \emph{bubble boundary} $\partial B_g$ is the set of marginal realizations at which
arbitrarily small perturbations may change the recognized symbol.
\end{definition}

The recognizability functional is defined as
\[
R(x) = \Pr\bigl[\pi_{\mathrm{glyph}}(x) = \text{intended glyph}\bigr],
\]
where the probability is over the relevant ensemble of observers and perceptual
conditions. Then $B_g = \{x : R(x) > R_{\min}\}$ for a threshold $R_{\min}$, and
$\partial B_g = \{x : R(x) = R_{\min}\}$. A realization $x$ with $R(x)$ well above
$R_{\min}$ is deeply within the admissibility bubble: it is robust to substantial
geometric perturbation. A realization with $R(x)$ near $R_{\min}$ is at the bubble
boundary: marginal perturbations may collapse it to a different projection class or
to no recognized symbol at all.

\subsection{Semantic Liftability Under Geometric Distortion}

The obstruction to admissible liftability proved in Part~IX applies directly to the
glyph system. An observer restricted to $M_{\mathrm{semantic}}$ --- an observer who
knows only the recognized symbolic identity, not the geometric realization --- cannot
reconstruct the specific element of $X_{\mathrm{stroke}}$ that produced it. This
is not an empirical limitation but a structural one. Many distinct geometric realizations
produce the same recognized symbol, and without additional information outside
$\mathcal{O}(M_{\mathrm{semantic}})$, no selection criterion distinguishes them.

The liftability obstruction for the glyph system can be stated as follows. The
structural symmetry group $G_{\pi_{\mathrm{glyph}}}$ of the glyph projection consists
of all geometric transformations of $X_{\mathrm{stroke}}$ that preserve the recognized
identity: stroke-width rescaling, compatible weight changes, admissible curvature
deformations, and rendering-medium variations that leave $\pi_{\mathrm{glyph}}$
invariant. This group acts freely on the fibers of $\pi_{\mathrm{glyph}}$, and the
associated Čech cohomology class $[\pi_{\mathrm{glyph}}] \in \check{H}^1(M_{\mathrm{semantic}},
\underline{G_{\pi_{\mathrm{glyph}}}})$ is non-trivial, preventing any admissible section.

This non-triviality is the formal statement that no canonical font exists: there is no
privileged geometric realization of any glyph that could serve as the unique
representative of its projection class without making an arbitrary choice of
fiber element.

\subsection{The RSVP Dynamics of Font Transformation}

Font transformations — from regular to bold, from upright to italic, from formal to
frenetic, from legible to hypnagogic — are flows on $X_{\mathrm{stroke}}$ parameterized
by the RSVP field triple $(\Phi, \mathbf{v}, S)$. In this context:

The scalar field $\Phi : X_{\mathrm{stroke}} \to \mathbb{R}$ measures the
\emph{recognizability density}: $\Phi(x)$ is high when $x$ is deep within the
admissibility bubble $B_g$ and decreases toward zero as $x$ approaches $\partial B_g$.
High $\Phi$ indicates a realization far from collapse; low $\Phi$ indicates proximity
to the bubble boundary.

The vector field $\mathbf{v} : X_{\mathrm{stroke}} \to TX_{\mathrm{stroke}}$ governs
the \emph{stylistic deformation flow}: the direction and rate at which a geometric
realization moves through the trajectory space under a given typographic transformation.
Bold transformation is a flow in the direction of increasing stroke mass. Italic
transformation is a directional shear flow. Condensation is a compressive flow in
the horizontal direction. Each typographic operation defines a specific vector field
on $X_{\mathrm{stroke}}$.

The entropy field $S : X_{\mathrm{stroke}} \to \mathbb{R}$ measures the \emph{local
geometric disorder} of the realization: high $S$ corresponds to high local variability
in stroke geometry, many competing curvature interpretations, or high spatial frequency
noise. Low $S$ corresponds to clean, unambiguous geometric specification.

The admissibility conservation equation $\partial_t \Phi + \nabla \cdot (\Phi\,\mathbf{v})
= -\lambda S$ then expresses the constraint that recognizability density is transported
by the stylistic flow $\mathbf{v}$ and dissipated at a rate proportional to local
geometric entropy. A typographic transformation that increases $S$ near the bubble
boundary is a transformation that threatens recognizability.

\subsection{Bold and Italic as Low-Entropy Manifold Flows}

Bold and italic are the canonical examples of \emph{admissible low-entropy flows}:
transformations that move a glyph realization through $X_{\mathrm{stroke}}$ while
remaining well within the admissibility bubble $B_g$. They increase or redirect
geometric mass without crossing the bubble boundary.

A bold transformation $\phi_{\mathrm{bold}} : X_{\mathrm{stroke}} \to X_{\mathrm{stroke}}$
expands the stroke width uniformly, increasing the admissibility density $\Phi$ in the
sense of making the realization more visually prominent and reducing the probability of
misrecognition due to small-scale rendering artifacts. Formally, bold transformation
is a flow in the direction $\nabla\Phi_{\mathrm{mass}}$, the gradient of the stroke-mass
density component of $\Phi$. The entropy field $S$ changes minimally under bold
transformation in the interior of $B_g$: the geometric structure of the glyph is
preserved, only its weight is increased. The admissibility conservation equation is
satisfied with $\lambda S$ remaining small.

An italic transformation $\phi_{\mathrm{italic}} : x \mapsto Ax$ is a shear map
$A : X_{\mathrm{stroke}} \to X_{\mathrm{stroke}}$ with $A = I + \theta E_{12}$ for
a shear angle $\theta$ and the elementary matrix $E_{12}$ acting on the horizontal
displacement of vertical strokes. The shear introduces a directional vector flow:
$\mathbf{v}_{\mathrm{italic}} = \theta\,\partial_x$, a constant horizontal displacement
field. For small $\theta$ (typical typographic italic angles are $7$--$15$ degrees),
the flow remains within $B_g$: the shear is admissible because horizontal displacement
preserves the topological skeleton of the glyph and the perceptual projection
$\pi_{\mathrm{glyph}}$ is invariant under these shears.

Both bold and italic are examples of fiber-preserving flows in the sense of the
Projection Robustness Principle: $\pi_{\mathrm{glyph}}(\phi(x)) = \pi_{\mathrm{glyph}}(x)$
for all $x$ in the interior of $B_g$.

\subsection{Frenetic Typography and Tangential Entropy Injection}

Frenetic typography represents a qualitatively different flow regime: one in which
the entropy field $S$ is substantially elevated, but the entropy flow remains
predominantly tangential to $M_{\mathrm{semantic}}$. A frenetic realization of a glyph
introduces high local variability --- jagged edges, varying stroke weights, disrupted
continuity, irregular spacing --- while preserving enough of the topological skeleton
that the projection $\pi_{\mathrm{glyph}}$ remains defined and stable.

In RSVP terms, frenetic typography increases $S$ substantially but keeps
$(\nabla S)_\perp \approx 0$: the entropy gradient is approximately tangential to
$M_{\mathrm{semantic}}$ at the projected point. By the Bubble Boundary as Entropy
Tangency Threshold proposition, the realization therefore remains within $B_g$ despite
elevated entropy. The recognizability density $\Phi$ decreases under the entropy
injection --- the realization is less far from the bubble boundary than a clean regular
rendering --- but does not reach zero.

This explains the phenomenology of frenetic typography: it reads as energetic,
turbulent, or chaotic, because the elevated $S$ is perceptible as visual noise, but it
remains readable because the projection $\pi_{\mathrm{glyph}}$ is still well-defined.
The fiber of the projection has expanded under entropy injection (more diverse
geometric realizations map to the same glyph), making the system less fragile in
some directions while also reducing the distance to $\partial B_g$.

Formally, define the frenetic flow as
\[
\phi_{\mathrm{fren}} : x \mapsto x + \xi(x),
\]
where $\xi : X_{\mathrm{stroke}} \to X_{\mathrm{stroke}}$ is a high-frequency
perturbation field satisfying $\mathbf{E}[\xi] = 0$ (zero mean) and
$\mathrm{Var}[\xi]$ large but confined to directions tangential to $M_{\mathrm{semantic}}$.
The condition $\mathrm{Var}[\xi]_\perp \approx 0$ is the entropy tangency condition;
it is the mathematical content of "frenetic but readable."

\subsection{Hypnagogic Typography and Boundary Oscillation}

Hypnagogic typography represents the most geometrically delicate regime: that of
slow oscillation near the bubble boundary $\partial B_g$. A hypnagogic realization
does not inject high-frequency entropy (as frenetic does) but instead introduces
low-frequency drift in the admissibility field, weakening the local curvature constraints
that define the glyph's geometric identity and allowing the projection class to become
momentarily ambiguous.

In RSVP terms, hypnagogic typography produces a slowly varying admissibility field
$\Phi(x, t)$ that oscillates near $\Phi_{\min}$: the system spends time near $\partial B_g$
without permanently crossing it. The entropy field $S$ is not elevated in total, but
its transverse component $(\nabla S)_\perp$ periodically becomes non-zero before
relaxing back to zero. The realization oscillates between stable projection (within
$B_g$) and marginal projection (on $\partial B_g$), never committing fully to
catastrophic collapse.

This explains the perceptual phenomenology of hypnagogic typography: it feels dreamlike
or liminal because the observer's perceptual projection $\pi_{\mathrm{glyph}}$ is
periodically destabilized. The glyph is recognized, then partially dissolved, then
re-recognized, in a cycle governed by the slow oscillation of $\Phi(x,t)$ around
$\Phi_{\min}$. Unlike frenetic typography, which is turbulent but stable, hypnagogic
typography is smooth but marginally stable. Unlike catastrophic collapse, which
permanently exits $B_g$, hypnagogic typography returns to $B_g$ after each excursion.

The formal model is a trajectory $x(t) \in X_{\mathrm{stroke}}$ satisfying the
damped oscillation
\[
\ddot{\Phi}(x(t)) + \gamma\,\dot{\Phi}(x(t)) + \omega_0^2\,(\Phi(x(t)) - \Phi_{\min}) = \eta(t),
\]
where $\gamma$ is a damping coefficient, $\omega_0$ is the natural oscillation frequency
of the admissibility field around $\Phi_{\min}$, and $\eta(t)$ is a low-frequency
noise term representing the slow geometric drift. The realization spends most of its
time near $\Phi_{\min}$, occasionally dipping below (brief collapse, perceptual
ambiguity) and recovering (return to recognition).

\subsection{Recognizability Thresholds and Catastrophic Collapse}

Catastrophic collapse in typography is the crossing of $\partial B_g$ without return:
a deformation so large that $\pi_{\mathrm{glyph}}(x)$ no longer maps to the intended
glyph, and the observer either recognizes a different symbol or recognizes nothing.
This is the typographic instance of the failure mode formalized in Part~VII.

The condition for catastrophic collapse is
\[
\exists\, \delta x : \pi_{\mathrm{glyph}}(x + \delta x) \neq g
\quad \text{and} \quad
\pi_{\mathrm{glyph}}(x + \delta x) \notin B_g.
\]
The realization has not merely moved to a different fiber of the same glyph but has
exited the admissibility bubble entirely, entering either the fiber of a different
glyph (a misrecognition) or the complement of all admissible regions (illegibility).

Recognizability as a function of deformation strength follows a characteristic
profile. For small deformations within $B_g$, $R(x+\delta x) \approx R(x) > R_{\min}$:
the glyph is robustly recognized. As the deformation approaches $\partial B_g$,
$R(x+\delta x)$ decreases toward $R_{\min}$: recognition becomes uncertain. At
$\partial B_g$, $R(x) = R_{\min}$: recognition is marginal. Beyond $\partial B_g$,
$R$ drops sharply as the perceptual system's projection $\pi_{\mathrm{glyph}}$ either
misidentifies the symbol or fails entirely.

The transition from marginal to catastrophic is a phase transition in the RSVP sense:
the Lyapunov function $\mathcal{V}(x) = V(\Phi(x)) + \alpha S(x)$ passes through a
critical point at $\partial B_g$, and the trajectory exits the basin of attraction of
$M_{\mathrm{semantic}}$ at $g$. The sharp drop in $R$ at the boundary corresponds
to the exit from the attractor basin proved in the Dynamical Attractor Theorem.

\subsection{Reading as Projection Recovery}

The typographic section makes a claim that goes beyond typography: reading itself is
a form of admissible projection recovery. An observer reading a text does not reconstruct
the exact geometric parameters of each glyph. The observer computes
$\pi_{\mathrm{glyph}}(x)$ from noisy, partial, perspective-distorted, illumination-variable
sensory data $x$, extracting the projection class without reconstructing the preimage.

This is the dynamical attractor theorem applied to perceptual cognition. The observer's
neural computation is an approximation to the RSVP flow that drives trajectory
representations toward the attractor basin of $M_{\mathrm{semantic}}$. The perceptual
system is robust to optical distortion, font variation, partial occlusion, and rendering
noise because all of these are perturbations within the admissibility bubble: they
change $x$ while preserving $\pi_{\mathrm{glyph}}(x)$.

Dyslexia, under this framework, is a condition in which the effective $\partial B_g$
for certain glyph pairs is pathologically narrow: the admissibility bubbles of visually
similar symbols overlap or are insufficiently separated, making the projection
$\pi_{\mathrm{glyph}}$ unstable near the boundaries between similar glyphs such as
b/d or p/q. The observer's perceptual system cannot reliably determine which basin of
attraction the sensory input falls into.

Adversarial typography --- the deliberate construction of glyphs that fool optical
character recognition systems while remaining human-readable --- is the construction
of realizations $x$ near $\partial B_g^{\mathrm{human}}$ but far inside
$B_g^{\mathrm{OCR}}$ for a different glyph, exploiting the fact that human and
machine perceptual systems define different admissibility bubbles over the same
trajectory space. The adversarial glyph sits at the intersection of two different
bubble boundaries, recognized as $g$ by one projection system and as $g'$ by another.

\subsection{Toward Admissible Generative Typography}

The framework developed in this part suggests a principled generative typography system
distinct from existing approaches. Current generative and style-transfer systems
optimize for statistical similarity to a training distribution, using objectives such
as perceptual loss, feature matching, or discriminator feedback. These objectives do
not explicitly model the admissibility bubble structure: a system optimizing perceptual
loss can produce outputs near $\partial B_g$ or even outside it without penalty, as
long as the statistical features match.

An \emph{admissible generative typography} system would instead explicitly maintain
the constraint $\Phi(x) > \Phi_{\min}$ throughout generation, ensuring that every
generated realization remains within the admissibility bubble. The RSVP field triple
provides the architecture: the scalar field $\Phi$ serves as a learned recognizability
density estimator; the vector field $\mathbf{v}$ defines the stylistic deformation
flow for each typographic operation; the entropy field $S$ controls the degree and
direction of geometric disorder introduced at each step.

Define the generation process as a constrained variational problem:
\[
\phi^* = \operatorname{argmax}_{\phi \in \mathcal{F}} \int_{X_{\mathrm{stroke}}}
\bigl[\mathcal{L}_{\mathrm{style}}(x) - \mu\,\mathbf{1}[\Phi(\phi(x)) < \Phi_{\min}]\bigr]
\,d\mu(x),
\]
where $\mathcal{L}_{\mathrm{style}}$ is the stylistic objective, $\mu > 0$ is a penalty
weight for admissibility violations, and the constraint $\Phi(\phi(x)) \geq \Phi_{\min}$
enforces that generated realizations remain within the admissibility bubble.

The four typographic regimes --- regular, bold, frenetic, hypnagogic --- correspond
to four different constraint geometries on this optimization. Regular maintains low $S$
throughout; bold increases $\Phi$ by expanding stroke mass; frenetic maximizes
$\mathrm{Var}[\xi]$ subject to the entropy tangency constraint
$(\nabla S)_\perp \approx 0$; hypnagogic maximizes the time spent near $\Phi_{\min}$
subject to the constraint that $\Phi$ oscillates above $\Phi_{\min}$ rather than
falling below it.

The theoretical contribution of this framework is: it separates \emph{semantic
stability} from \emph{geometric specificity}. A font can be arbitrarily geometrically
novel --- unprecedented stroke geometry, never-before-seen weight distribution, entirely
original spatial frequency profile --- while remaining semantically stable because it
remains within the admissibility bubble. Conversely, a font that closely resembles
a standard typeface in geometric parameters can be semantically unstable if its
realizations concentrate near $\partial B_g$.

This separation is what allows function to survive collapse: the semantic invariant
$\pi_{\mathrm{glyph}}(x)$ is preserved not by preserving geometric fidelity but by
remaining within the correct basin of attraction in the RSVP field geometry.

%------------------------------------------------------
\section{Conclusion: What Survives Collapse}
%------------------------------------------------------

Three results from distinct scientific domains have been examined in this essay.
The first shows that anonymous additive aggregation suffices for universal secure
computation, because local entropy injection can destroy provenance while preserving
functional output. The second shows that bacterial gene expression is not a sequential
pipeline but a geometry-constrained coordination of partially coupled flows, where mRNA
fate depends on spatial accessibility rather than informational sequence. The third
shows that early seed development in \textit{Arabidopsis} is not a centrally orchestrated
blueprint execution but a distributed constraint-propagation system in which symplastically
isolated tissue compartments maintain developmental coherence through low-dimensional
coordination operators across admissibility interfaces.

All three results displace the pipeline ontology from their respective domains.
In computation, the pipeline from sender to receiver is replaced by local randomization
plus additive aggregation. In prokaryotic gene regulation, the pipeline from DNA to
protein is replaced by a partially coupled dynamical system in which spatial geometry,
kinetic accessibility, and entropy management collectively determine which observables
are stable. In plant development, the pipeline from genome to organism is replaced by
a distributed field system in which tissue-autonomous trajectories are coupled through
sparse low-bandwidth signalling projections across constrained interfaces.

The admissible projection framework formalizes what all three results share: a map
$\pi : X \to M$ that preserves computable functional observables while destroying
microscopic recoverability. The mechanisms differ --- mixing in the cryptographic case,
topological separation in the bacterial case, sparse inter-manifold coordination in the
developmental case --- but the abstract structure is the same in each. None of the
three systems preserves microscopic trajectory information. All three maintain functional
invariants. They are realizations, through different physical substrates and at
different scales, of the impossibility of local-to-global reconstruction in a
constrained topological setting and of the persistence of function despite that
impossibility.

The RSVP field triple $(\Phi, \mathbf{v}, S)$ provides the overarching dynamical
geometry, with the admissibility conservation equation $\partial_t \Phi + \nabla \cdot
(\Phi\,\mathbf{v}) = -\lambda S$ expressing the fundamental tension between trajectory
flow and entropy production. Functional persistence requires that entropy flow remain
tangent to the admissible manifold. When it does, entropy is a resource that enlarges
fibers and deepens robustness. When it does not, entropy is a perturbation that drives
trajectories out of the admissible regime and collapses function.

The deeper implication is a reorientation of scientific ontology. Classical systems
theory asks how trajectories are preserved. The admissible projection framework asks
which invariants survive trajectory destruction. That inversion is not a reformulation.
It is a recognition that the systems most worth understanding --- biological regulation,
secure computation, developmental coordination, cognitive abstraction, morphogenesis ---
are precisely those in which trajectories are not preserved, and yet function persists
anyway. Stability belongs not to the microscopic realization but to the projection
class itself.

The microscopic paths are irretrievable. The functional output is stable. The projection
is irreversible. The invariant persists.

The pipeline breaks. The invariant holds.

%------------------------------------------------------
\newpage
\bibliographystyle{plainnat}
\begin{thebibliography}{9}

\bibitem[Halevi et al.(2023)]{halevi2023additive}
Halevi, S., Ishai, Y., Kushilevitz, E., and Rabin, T.\ (2023).
Additive randomized encodings and their applications.
Cryptology ePrint Archive, Report 2023/870.

\bibitem[Kim et al.(2026)]{kim2026spatial}
Kim, S., Zhang, Y., Ju, X., Hassan, A., Wang, Y.-H., Liu, S., and Kim, S.\ (2026).
Spatial and genetic constraints govern transcription--translation coupling
and mRNA degradation in bacteria.
\textit{Nature Microbiology}.
\url{https://doi.org/10.1038/s41564-026-02374-8}

\bibitem[Bitansky et al.(2025)]{shufflinguniversal}
Bitansky, N., Erabelli, S., Garg, R., and Ishai, Y.\ (2025).
Shuffling is universal: Statistical additive randomized encodings for all functions.
Cryptology ePrint Archive, Paper 2025/1442.
\url{https://eprint.iacr.org/2025/1442}

\bibitem[Martin et al.(2026)]{martin2026atlas}
Martin, C.A., Cogdill, K.R., Pusey, A.L., et al.\ (2026).
A transcriptional atlas of early \textit{Arabidopsis} seed development suggests
mechanisms for inter-tissue coordination.
\textit{Nature Plants} \textbf{12}, 1133--1149.
\url{https://doi.org/10.1038/s41477-026-02295-8}

\end{thebibliography}

\end{document}
